% ============================================================
%  FURTHER PROBABILITY & STATISTICS
%  Comprehensive University-Preparation Notes
%  Author: Rayyan Hanif, Kolej MARA Seremban
% ============================================================

\documentclass[12pt, a4paper, oneside]{book}

% ── Core packages ──────────────────────────────────────────
\usepackage[T1]{fontenc}
\usepackage[utf8]{inputenc}
\usepackage{microtype}
\usepackage[a4paper, top=2.5cm, bottom=2.5cm, inner=3cm, outer=2.5cm]{geometry}

% ── Mathematics ────────────────────────────────────────────
\usepackage{amsmath}
\usepackage{amssymb}
\usepackage{amsthm}
\usepackage{mathtools}
\usepackage{bm}
\usepackage{siunitx}
\pdfprotrudechars=0
\pdfadjustspacing=0
\usepackage[T1]{fontenc}
\usepackage{lmodern}

% ── Graphics ───────────────────────────────────────────────
\usepackage{tikz}
\usepackage{pgfplots}
\pgfplotsset{compat=1.18}
\usetikzlibrary{arrows.meta, calc, patterns, shapes.geometric,
                decorations.pathmorphing, decorations.markings,
                positioning, fit, backgrounds, matrix, trees}
\usepgfplotslibrary{fillbetween, statistics}

% ── Colour ─────────────────────────────────────────────────
\usepackage[dvipsnames, table]{xcolor}

% ── Boxes ──────────────────────────────────────────────────
\usepackage[most]{tcolorbox}
\tcbuselibrary{theorems, breakable, skins}

% ── Tables ─────────────────────────────────────────────────
\usepackage{booktabs}
\usepackage{array}
\usepackage{multirow}
\usepackage{longtable}

% ── Lists ──────────────────────────────────────────────────
\usepackage{enumitem}

% ── Typography & Layout ────────────────────────────────────
\usepackage{parskip}
\usepackage{setspace}
\usepackage{fancyhdr}
\usepackage{titlesec}
\usepackage{emptypage}


% ── Hyperlinks ─────────────────────────────────────────────
\usepackage[hidelinks, colorlinks=false]{hyperref}

% ── Index ──────────────────────────────────────────────────
\usepackage{makeidx}
\makeindex

% ══════════════════════════════════════════════════════════
%  COLOUR PALETTE
% ══════════════════════════════════════════════════════════
\definecolor{PrimaryBlue}{RGB}{0, 84, 166}
\definecolor{AccentTeal}{RGB}{0, 150, 136}
\definecolor{AccentOrange}{RGB}{230, 100, 0}
\definecolor{AccentRed}{RGB}{198, 40, 40}
\definecolor{AccentGreen}{RGB}{27, 120, 55}
\definecolor{LightBlue}{RGB}{220, 235, 255}
\definecolor{LightTeal}{RGB}{200, 240, 235}
\definecolor{LightOrange}{RGB}{255, 235, 210}
\definecolor{LightRed}{RGB}{255, 220, 220}
\definecolor{LightGreen}{RGB}{215, 245, 225}
\definecolor{LightGrey}{RGB}{248, 248, 250}
\definecolor{DarkGrey}{RGB}{60, 60, 70}
\definecolor{MidGrey}{RGB}{140, 140, 155}

% ══════════════════════════════════════════════════════════
%  THEOREM ENVIRONMENTS
% ══════════════════════════════════════════════════════════
\theoremstyle{definition}
\newtheorem{definition}{Definition}[section]
\newtheorem{example}{Example}[section]
\newtheorem{exercise}{Exercise}[section]

\theoremstyle{plain}
\newtheorem{theorem}{Theorem}[section]
\newtheorem{proposition}{Proposition}[section]
\newtheorem{lemma}{Lemma}[section]
\newtheorem{corollary}{Corollary}[theorem]

\theoremstyle{remark}
\newtheorem{remark}{Remark}[section]
\newtheorem{note}{Note}[section]

% ══════════════════════════════════════════════════════════
%  TCOLORBOX STYLES
% ══════════════════════════════════════════════════════════

% Intuition box
\tcbset{intuitionstyle/.style={
  enhanced, breakable,
  colback=LightBlue, colframe=PrimaryBlue,
  boxrule=0.8pt, arc=4pt,
  title={\faLightbulb\quad Intuition},
  fonttitle=\bfseries\color{white},
  coltitle=white,
  attach boxed title to top left={yshift=-2mm, xshift=4mm},
  boxed title style={colback=PrimaryBlue, arc=3pt, boxrule=0pt},
  top=6mm
}}

\newtcolorbox{intuition}[1][]{
  enhanced, breakable,
  colback=LightBlue, colframe=PrimaryBlue,
  boxrule=0.8pt, arc=4pt,
  title={\textbf{Intuition} #1},
  fonttitle=\bfseries\color{white},
  coltitle=white,
  attach boxed title to top left={yshift=-2mm, xshift=4mm},
  boxed title style={colback=PrimaryBlue, arc=3pt, boxrule=0pt},
  top=6mm
}

% Warning box
\newtcolorbox{warningbox}[1][]{
  enhanced, breakable,
  colback=LightRed, colframe=AccentRed,
  boxrule=0.8pt, arc=4pt,
  title={\textbf{Warning} #1},
  fonttitle=\bfseries\color{white},
  coltitle=white,
  attach boxed title to top left={yshift=-2mm, xshift=4mm},
  boxed title style={colback=AccentRed, arc=3pt, boxrule=0pt},
  top=6mm
}

% Exam trap box
\newtcolorbox{examtrap}[1][]{
  enhanced, breakable,
  colback=LightOrange, colframe=AccentOrange,
  boxrule=0.8pt, arc=4pt,
  title={\textbf{Exam Trap} #1},
  fonttitle=\bfseries\color{white},
  coltitle=white,
  attach boxed title to top left={yshift=-2mm, xshift=4mm},
  boxed title style={colback=AccentOrange, arc=3pt, boxrule=0pt},
  top=6mm
}

% Historical note
\newtcolorbox{historicalbox}[1][]{
  enhanced, breakable,
  colback=LightGrey, colframe=MidGrey,
  boxrule=0.8pt, arc=4pt,
  title={\textbf{Historical Note} #1},
  fonttitle=\bfseries\color{DarkGrey},
  coltitle=DarkGrey,
  attach boxed title to top left={yshift=-2mm, xshift=4mm},
  boxed title style={colback=MidGrey, arc=3pt, boxrule=0pt},
  top=6mm
}

% Real-world application
\newtcolorbox{realworld}[1][]{
  enhanced, breakable,
  colback=LightGreen, colframe=AccentGreen,
  boxrule=0.8pt, arc=4pt,
  title={\textbf{Real-World Application} #1},
  fonttitle=\bfseries\color{white},
  coltitle=white,
  attach boxed title to top left={yshift=-2mm, xshift=4mm},
  boxed title style={colback=AccentGreen, arc=3pt, boxrule=0pt},
  top=6mm
}

% Exploration box
\newtcolorbox{exploration}[1][]{
  enhanced, breakable,
  colback=LightTeal, colframe=AccentTeal,
  boxrule=0.8pt, arc=4pt,
  title={\textbf{Exploration} #1},
  fonttitle=\bfseries\color{white},
  coltitle=white,
  attach boxed title to top left={yshift=-2mm, xshift=4mm},
  boxed title style={colback=AccentTeal, arc=3pt, boxrule=0pt},
  top=6mm
}

% Key result box
\newtcolorbox{keyresult}[1][]{
  enhanced, breakable,
  colback=LightBlue!60!white, colframe=PrimaryBlue!80!black,
  boxrule=1.2pt, arc=5pt,
  title={\textbf{Key Result} #1},
  fonttitle=\bfseries\color{white},
  coltitle=white,
  attach boxed title to top left={yshift=-2mm, xshift=4mm},
  boxed title style={colback=PrimaryBlue!80!black, arc=3pt, boxrule=0pt},
  top=6mm
}

% Worked example box
\newtcolorbox[auto counter, number within=section]{workedexample}[1][]{
  enhanced, breakable,
  colback=white, colframe=AccentTeal,
  boxrule=0.6pt, arc=4pt,
  title={\textbf{Worked Example \thetcbcounter} #1},
  fonttitle=\bfseries\color{white},
  coltitle=white,
  attach boxed title to top left={yshift=-2mm, xshift=4mm},
  boxed title style={colback=AccentTeal, arc=3pt, boxrule=0pt},
  top=6mm
}

% Summary box
\newtcolorbox{summarybox}[1][]{
  enhanced, breakable,
  colback=PrimaryBlue!8!white, colframe=PrimaryBlue,
  boxrule=1pt, arc=5pt,
  title={\textbf{Summary} #1},
  fonttitle=\bfseries\color{white},
  coltitle=white,
  attach boxed title to top left={yshift=-2mm, xshift=4mm},
  boxed title style={colback=PrimaryBlue, arc=3pt, boxrule=0pt},
  top=6mm
}

% ══════════════════════════════════════════════════════════
%  HEADER / FOOTER
% ══════════════════════════════════════════════════════════
\pagestyle{fancy}
\fancyhf{}
\fancyhead[L]{\small\textcolor{MidGrey}{\leftmark}}
\fancyhead[R]{\small\textcolor{MidGrey}{Further Probability \& Statistics}}
\fancyfoot[C]{\small\textcolor{MidGrey}{\thepage}}
\renewcommand{\headrulewidth}{0.4pt}
\renewcommand{\footrulewidth}{0pt}

% ══════════════════════════════════════════════════════════
%  TITLE FORMATTING
% ══════════════════════════════════════════════════════════
\titleformat{\chapter}[display]
  {\normalfont\huge\bfseries\color{PrimaryBlue}}
  {\chaptertitlename\ \thechapter}{20pt}{\Huge}
\titlespacing*{\chapter}{0pt}{10pt}{30pt}

\titleformat{\section}
  {\normalfont\Large\bfseries\color{PrimaryBlue!80!black}}
  {\thesection}{1em}{}
\titleformat{\subsection}
  {\normalfont\large\bfseries\color{DarkGrey}}
  {\thesubsection}{1em}{}
\titleformat{\subsubsection}
  {\normalfont\normalsize\bfseries\itshape\color{DarkGrey}}
  {\thesubsubsection}{1em}{}

% ══════════════════════════════════════════════════════════
%  CUSTOM COMMANDS
% ══════════════════════════════════════════════════════════
\newcommand{\E}{\mathbb{E}}
\newcommand{\Var}{\operatorname{Var}}
\newcommand{\Cov}{\operatorname{Cov}}
\newcommand{\Corr}{\operatorname{Corr}}
\newcommand{\Prob}{\mathbb{P}}
\newcommand{\R}{\mathbb{R}}
\newcommand{\N}{\mathbb{N}}
\newcommand{\indep}{\perp\!\!\!\perp}
\newcommand{\iid}{\overset{\text{iid}}{\sim}}
\newcommand{\xbar}{\bar{X}}
\newcommand{\pvalue}{p\text{-value}}
\newcommand{\df}{\nu}
\newcommand{\chistat}{\chi^2}
\DeclareMathOperator{\sign}{sign}
\DeclareMathOperator{\rank}{rank}

% Solution environment
\newenvironment{solution}
  {\par\medskip\noindent\textbf{Solution.}\quad}
  {\hfill$\square$\par\medskip}

% ══════════════════════════════════════════════════════════
%  DOCUMENT BEGINS
% ══════════════════════════════════════════════════════════
\begin{document}

% ── Title Page ─────────────────────────────────────────────
\begin{titlepage}
\begin{tikzpicture}[remember picture, overlay]
  % Background strip
  \fill[PrimaryBlue] (current page.north west) rectangle
        ([yshift=-0.38\paperheight]current page.north east);
  % Decorative arc
  \fill[AccentTeal!60!PrimaryBlue, opacity=0.6]
        ([yshift=-0.38\paperheight]current page.north west)
        .. controls +(4,0) and +(-4,0) ..
        ([yshift=-0.30\paperheight]current page.north east)
        -- (current page.north east) -- (current page.north west) -- cycle;
  % Bottom strip
  \fill[PrimaryBlue!20!white] (current page.south west) rectangle
        ([yshift=1.8cm]current page.south east);
\end{tikzpicture}

\vspace*{3.5cm}
{\color{white}\fontsize{32}{40}\selectfont\bfseries Further Probability\\[6pt] \& Statistics\par}
\vspace{0.8cm}
{\color{white!80!PrimaryBlue}\fontsize{16}{22}\selectfont
 University-Preparation Notes\\[4pt]
 Chapters 8--12\par}

\vspace{3cm}
\noindent
{\fontsize{13}{18}\selectfont\color{DarkGrey}
\begin{tabular}{@{}ll}
\textbf{Author:}  & Rayyan Hanif \\[4pt]
\textbf{Institution:} & Kolej MARA Seremban \\[4pt]
\textbf{Level:}   & A-Level Further Mathematics \\[4pt]
\textbf{Audience:} & Aspiring Cambridge, Oxford, Imperial, \\
                   & Warwick, MIT, Stanford applicants \\
\end{tabular}}

\vfill
\noindent\rule{\linewidth}{0.4pt}\\[4pt]
{\small\color{MidGrey}
These notes build deep mathematical understanding of Further Statistics.
Topics include continuous random variables, inferential statistics,
chi-squared tests, non-parametric tests, and probability generating functions.}
\end{titlepage}

% ── Preface ────────────────────────────────────────────────
\chapter*{Preface}
\addcontentsline{toc}{chapter}{Preface}

These notes are written for the ambitious Further Mathematics student whose
sights are set beyond merely passing examinations.  The goal is to build
\emph{genuine statistical thinking}: the ability to reason carefully about
uncertainty, to choose appropriate methods, to understand what those methods
actually do, and to interpret results correctly.

Statistics is not a collection of procedures to be memorised.  It is a
principled framework for learning from data in the presence of randomness.
Every formula here has a story behind it---a problem someone needed to solve,
an insight someone had, and a derivation that makes the formula inevitable
rather than arbitrary.

\medskip
\textbf{How to use these notes.}\quad
Read each section in order the first time.  When you encounter a result,
try to anticipate the derivation before reading it.  Work every example
yourself before looking at the solution.  The exploration and extension
questions at the end of each chapter are designed to take you beyond the
syllabus into the territory you will encounter in your first year of
university.

\medskip
\textbf{Prerequisites.}\quad
These notes assume familiarity with A-Level Mathematics (probability,
statistics, calculus) and the earlier chapters of Further Statistics
(discrete distributions, sampling, hypothesis testing with the normal
and Poisson distributions, linear regression and correlation).

\tableofcontents

% ══════════════════════════════════════════════════════════
%  CHAPTER 8: CONTINUOUS RANDOM VARIABLES
% ══════════════════════════════════════════════════════════
\chapter{Continuous Random Variables}
\label{ch:crv}

\section{Introduction}

So far you have studied \emph{discrete} random variables: quantities that
take values in a finite or countably infinite set (such as the number of
heads in ten coin flips, or the number of calls arriving at a switchboard
in one minute).  For discrete variables, we can assign a probability to
each individual outcome.

But many phenomena in the real world are \emph{inherently continuous}: the
exact height of a randomly selected adult, the time until a radioactive atom
decays, the return on a financial investment over a year.  For such
quantities, the outcome can be \emph{any} real number in some interval, and
it becomes meaningless to ask for the probability of any \emph{exact} value.

\begin{intuition}
If you pick a real number uniformly at random from $[0,1]$, what is the
probability that you pick exactly $0.5$?  The answer is zero---not because
it is impossible, but because the interval contains infinitely many points
and no single point has positive measure.  What \emph{does} make sense is
asking for the probability of landing in an interval, e.g.,
$\Prob(0.3 \le X \le 0.7) = 0.4$.
\end{intuition}

Continuous random variables arise throughout science and engineering:
\begin{itemize}[leftmargin=2em]
  \item \textbf{Physics:} lifetimes of particles, measurement errors.
  \item \textbf{Finance:} stock prices, option payoffs, interest rates.
  \item \textbf{Biology:} heights, weights, enzyme concentrations.
  \item \textbf{Engineering:} stress tolerances, signal noise, waiting times.
  \item \textbf{Machine learning:} model parameters treated as continuous
        latent variables.
\end{itemize}

\section{The Probability Density Function}
\label{sec:pdf}

\subsection{Definition and Intuition}

\begin{definition}[Probability Density Function]
A random variable $X$ is said to be \emph{continuous} if there exists a
non-negative function $f:\R\to[0,\infty)$, called the
\textbf{probability density function} (PDF) of $X$, such that for all
$a \le b$,
\[
  \Prob(a \le X \le b) \;=\; \int_a^b f(x)\,\mathrm{d}x.
\]
\end{definition}

The PDF must satisfy two conditions:
\begin{enumerate}[label=\textbf{(\roman*)}]
  \item \textbf{Non-negativity:} $f(x) \ge 0$ for all $x \in \R$.
  \item \textbf{Normalisation:} $\displaystyle\int_{-\infty}^{\infty} f(x)\,\mathrm{d}x = 1$.
\end{enumerate}

\begin{intuition}
Think of the PDF as a \emph{density of probability per unit length}.  A
height $f(x)$ does \emph{not} represent a probability; indeed $f(x)$ can be
greater than 1.  The \emph{area} under the curve between two points
represents the probability of the variable falling between those points.  
This is exactly the way a histogram works: the area of each bar gives the
relative frequency (proportion) of observations in that interval.
\end{intuition}

\begin{figure}[h]
\centering
\begin{tikzpicture}
\begin{axis}[
  width=13cm, height=6cm,
  xlabel={$x$}, ylabel={$f(x)$},
  xmin=-0.5, xmax=5.5, ymin=0, ymax=0.55,
  axis lines=left,
  xtick={1, 2.5, 4},
  xticklabels={$a$, , $b$},
  ytick=\empty,
  title={Probability as area under the PDF},
  title style={font=\bfseries},
]
% Normal-ish density
\addplot[thick, PrimaryBlue, domain=-0.2:5.5, samples=200]
  {0.45*exp(-0.5*(x-2.5)^2/1.2)};
% Shaded region
\addplot[fill=AccentTeal, fill opacity=0.35, draw=none, domain=1:4, samples=200]
  {0.45*exp(-0.5*(x-2.5)^2/1.2)} \closedcycle;
% Vertical lines
\draw[dashed, AccentTeal] (axis cs:1,0) -- (axis cs:1,{0.45*exp(-0.5*(1-2.5)^2/1.2)});
\draw[dashed, AccentTeal] (axis cs:4,0) -- (axis cs:4,{0.45*exp(-0.5*(4-2.5)^2/1.2)});
% Label for area
\node[font=\small] at (axis cs:2.5, 0.18) {$\Prob(a \le X \le b)$};
\end{axis}
\end{tikzpicture}
\caption{The probability $\Prob(a\le X\le b)$ equals the shaded area
  under the PDF between $x=a$ and $x=b$.}
\label{fig:pdf-area}
\end{figure}

\subsection{Key Properties}

Because probabilities correspond to areas rather than heights, several
important properties follow immediately:

\begin{proposition}
For a continuous random variable $X$ with PDF $f$:
\begin{enumerate}[label=\textbf{(\roman*)}]
  \item $\Prob(X = a) = 0$ for every $a \in \R$.
  \item $\Prob(a < X < b) = \Prob(a \le X < b) = \Prob(a < X \le b) = \Prob(a \le X \le b)$.
  \item $\Prob(X \ge a) = \int_a^\infty f(x)\,\mathrm{d}x$.
  \item $\Prob(X \le a) + \Prob(X > a) = 1$.
\end{enumerate}
\end{proposition}

Property (i) is crucial: the probability of any \emph{specific} value is
zero.  This does \emph{not} mean that value cannot occur---it simply reflects
the nature of the continuous setting.

\begin{examtrap}
Many students write $\Prob(X = a) = f(a)$.  This is \textbf{wrong}.
$f(a)$ is a \emph{density}, not a probability.  The probability of any
single point is zero.  Only integrals (areas) give probabilities.
\end{examtrap}

\subsection{Finding Probabilities: Worked Examples}

\begin{workedexample}[--- Verifying a PDF and Computing Probabilities]
Let $X$ have PDF
\[
  f(x) = \begin{cases} kx(4-x) & 0 \le x \le 4 \\ 0 & \text{otherwise.} \end{cases}
\]
\begin{enumerate}[label=\textbf{(\alph*)}]
  \item Find the value of $k$.
  \item Find $\Prob(1 \le X \le 3)$.
  \item Find $\Prob(X > 2.5)$.
\end{enumerate}

\begin{solution}
\textbf{(a)}  The normalisation condition requires
$\int_{-\infty}^{\infty} f(x)\,\mathrm{d}x = 1$.  Since $f(x)=0$
outside $[0,4]$:
\[
  \int_0^4 kx(4-x)\,\mathrm{d}x = k\int_0^4 (4x - x^2)\,\mathrm{d}x
  = k\left[2x^2 - \frac{x^3}{3}\right]_0^4
  = k\!\left(32 - \frac{64}{3}\right) = k\cdot\frac{32}{3}.
\]
Setting this equal to 1: $k = \dfrac{3}{32}$.

\medskip
\textbf{(b)}
\[
  \Prob(1\le X\le 3) = \frac{3}{32}\int_1^3 x(4-x)\,\mathrm{d}x
  = \frac{3}{32}\left[2x^2 - \frac{x^3}{3}\right]_1^3
  = \frac{3}{32}\!\left[\left(18 - 9\right) - \left(2 - \tfrac{1}{3}\right)\right]
  = \frac{3}{32}\cdot\frac{20}{3} = \frac{20}{32} = \frac{5}{8}.
\]

\textbf{(c)}
\[
  \Prob(X > 2.5) = \frac{3}{32}\int_{2.5}^{4} x(4-x)\,\mathrm{d}x
  = \frac{3}{32}\left[2x^2 - \frac{x^3}{3}\right]_{2.5}^{4}.
\]
At $x=4$: $2(16)-\frac{64}{3} = 32 - \frac{64}{3} = \frac{32}{3}$.
At $x=2.5$: $2(6.25)-\frac{15.625}{3} = 12.5 - \frac{15.625}{3} = \frac{22.375}{3}$.
\[
  \Prob(X>2.5) = \frac{3}{32}\cdot\frac{32-22.375}{3} = \frac{9.625}{32}
  = \frac{77}{256} \approx 0.301.
\]
\end{solution}
\end{workedexample}

\subsection{Common Continuous Distributions}

\begin{definition}[Uniform Distribution]
$X \sim U(a,b)$ if
\[
  f(x) = \frac{1}{b-a}, \quad a \le x \le b.
\]
Every sub-interval of equal length has the same probability.
\end{definition}

\begin{definition}[Exponential Distribution]
$X \sim \operatorname{Exp}(\lambda)$ if
\[
  f(x) = \lambda e^{-\lambda x}, \quad x \ge 0,\quad \lambda > 0.
\]
It models the time until the first event in a Poisson process with rate
$\lambda$.  It has the remarkable \emph{memoryless property}:
$\Prob(X > s+t \mid X > s) = \Prob(X > t)$.
\end{definition}

\begin{figure}[h]
\centering
\begin{tikzpicture}
\begin{axis}[
  width=13cm, height=5.5cm,
  xlabel={$x$}, ylabel={$f(x)$},
  xmin=0, xmax=5.2, ymin=0, ymax=1.15,
  axis lines=left,
  legend pos=north east,
  legend style={font=\small},
  title={Exponential PDFs for different rates $\lambda$},
  title style={font=\bfseries},
]
\addplot[thick, PrimaryBlue, domain=0:5.2, samples=200] {1.0*exp(-1.0*x)};
\addlegendentry{$\lambda=1$}
\addplot[thick, AccentTeal, domain=0:5.2, samples=200] {0.5*exp(-0.5*x)};
\addlegendentry{$\lambda=0.5$}
\addplot[thick, AccentOrange, domain=0:5.2, samples=200] {2.0*exp(-2.0*x)};
\addlegendentry{$\lambda=2$}
\end{axis}
\end{tikzpicture}
\caption{Exponential PDFs: larger $\lambda$ means faster decay and shorter
expected waiting times.}
\label{fig:exp-pdf}
\end{figure}

\section{The Cumulative Distribution Function}
\label{sec:cdf}

\subsection{Definition}

\begin{definition}[Cumulative Distribution Function]
The \textbf{cumulative distribution function} (CDF) of a random variable $X$
is defined for all $x \in \R$ by
\[
  F(x) = \Prob(X \le x).
\]
For a continuous random variable,
\[
  F(x) = \int_{-\infty}^{x} f(t)\,\mathrm{d}t.
\]
\end{definition}

The CDF accumulates probability as $x$ increases.

\subsection{Properties of the CDF}

\begin{proposition}[Properties of the CDF]
For any CDF $F$:
\begin{enumerate}[label=\textbf{(\roman*)}]
  \item $0 \le F(x) \le 1$ for all $x$.
  \item $F$ is non-decreasing: if $a < b$ then $F(a) \le F(b)$.
  \item $\lim_{x\to-\infty} F(x) = 0$ and $\lim_{x\to+\infty} F(x) = 1$.
  \item For a \emph{continuous} RV, $F$ is continuous everywhere.
  \item The PDF is the derivative of the CDF:
        $f(x) = F'(x)$ wherever the derivative exists.
\end{enumerate}
\end{proposition}

The fifth property is the Fundamental Theorem of Calculus applied to
probability.  It tells us that the PDF and CDF are two sides of the same
coin: one is the derivative of the other; the other is the integral of the first.

\begin{figure}[h]
\centering
\begin{tikzpicture}
\begin{axis}[
  width=13cm, height=6cm,
  xlabel={$x$}, ylabel={},
  xmin=-0.5, xmax=5.5, ymin=0, ymax=1.1,
  axis lines=left,
  legend pos=south east,
  legend style={font=\small},
  title={PDF (area interpretation) and its CDF},
  title style={font=\bfseries},
]
% PDF scaled to fit
\addplot[thick, PrimaryBlue, domain=-0.2:5.5, samples=200]
  {0.9*exp(-0.5*(x-2.5)^2/1.2)};
\addlegendentry{PDF $f(x)$}
% CDF
\addplot[thick, AccentOrange, domain=-0.2:5.5, samples=200]
  {1/(1+exp(-2*(x-2.5)))};
\addlegendentry{CDF $F(x)$}
% Horizontal asymptote
\draw[dashed, MidGrey] (axis cs:-0.5,1) -- (axis cs:5.5,1)
  node[right, font=\small] {$1$};
\end{axis}
\end{tikzpicture}
\caption{The CDF starts at 0, is S-shaped and non-decreasing, and
  asymptotes to 1.  It is the running integral of the PDF.}
\label{fig:cdf}
\end{figure}

\subsection{Using the CDF}

Once we have $F(x)$, computing probabilities is straightforward:
\[
  \Prob(a \le X \le b) = F(b) - F(a), \qquad
  \Prob(X > a) = 1 - F(a).
\]

\begin{workedexample}[--- Deriving the CDF]
For the PDF $f(x) = \frac{3}{32}x(4-x)$, $0 \le x \le 4$, find $F(x)$ and
hence find the median.

\begin{solution}
For $x < 0$: $F(x)=0$.\\
For $0 \le x \le 4$:
\[
  F(x) = \int_0^x \frac{3}{32}t(4-t)\,\mathrm{d}t
  = \frac{3}{32}\left[2t^2 - \frac{t^3}{3}\right]_0^x
  = \frac{3}{32}\!\left(2x^2 - \frac{x^3}{3}\right)
  = \frac{x^2}{16}\!\left(3 - \frac{x}{4}\right).
\]
For $x > 4$: $F(x) = 1$.

\medskip
\textbf{Median $m$:} Solve $F(m) = 0.5$:
\[
  \frac{m^2(12 - m)}{64} = 0.5 \implies m^2(12-m) = 32.
\]
This is a cubic.  The distribution is symmetric about $x=2$ (check: the PDF
$f(x) = f(4-x)$), so the median is $m = 2$.
\end{solution}
\end{workedexample}

\subsection{Quantiles, Percentiles, and the Median}

\begin{definition}[Quantile]
The \textbf{$p$-th quantile} (or $100p$-th percentile) of a continuous random
variable $X$ is the value $x_p$ satisfying $F(x_p) = p$,
equivalently $\Prob(X \le x_p) = p$.
\end{definition}

Special cases:
\begin{itemize}
  \item $p = 0.5$: the \textbf{median} $m$.
  \item $p = 0.25$: the \textbf{lower quartile} $Q_1$.
  \item $p = 0.75$: the \textbf{upper quartile} $Q_3$.
\end{itemize}

For a symmetric unimodal distribution, mean = median = mode.

\section{Expectation of a Function of a Continuous Random Variable}
\label{sec:egx}

\subsection{Expected Value}

\begin{definition}[Expected Value]
For a continuous RV $X$ with PDF $f$, the \textbf{expected value}
(mean) is
\[
  \E(X) = \int_{-\infty}^{\infty} x\,f(x)\,\mathrm{d}x,
\]
provided the integral converges absolutely.
\end{definition}

\begin{intuition}
$\E(X)$ is the long-run average value of $X$ over many independent
repetitions.  It is the \emph{centre of mass} of the probability density:
if you placed mass $f(x)\,\mathrm{d}x$ at position $x$ along a weightless
rod, $\E(X)$ is where you would place a fulcrum to balance it.
\end{intuition}

\subsection{The Law of the Unconscious Statistician (LOTUS)}

\begin{theorem}[LOTUS for Continuous RVs]
\label{thm:lotus}
If $X$ has PDF $f$ and $g:\R\to\R$ is a measurable function, then
\[
  \E(g(X)) = \int_{-\infty}^{\infty} g(x)\,f(x)\,\mathrm{d}x,
\]
provided the integral converges absolutely.
\end{theorem}

\begin{remark}
The name is humorous: students sometimes forget this theorem and compute
$\E(g(X))$ by first finding the distribution of $g(X)$ and then integrating.
The theorem says you can use the original PDF---a major computational
shortcut.
\end{remark}

Key applications of LOTUS:
\begin{align*}
  \E(X^2) &= \int_{-\infty}^{\infty} x^2 f(x)\,\mathrm{d}x, \\
  \Var(X) &= \E(X^2) - [\E(X)]^2, \\
  \E(aX+b) &= a\E(X) + b.
\end{align*}

\begin{definition}[Variance and Standard Deviation]
\[
  \Var(X) = \E\bigl[(X-\mu)^2\bigr] = \E(X^2) - \mu^2,
  \quad \text{where } \mu = \E(X).
\]
The \textbf{standard deviation} is $\sigma = \sqrt{\Var(X)}$.
\end{definition}

\begin{workedexample}[--- Expectation and Variance]
For $f(x) = \frac{3}{32}x(4-x)$, $0\le x\le 4$, find $\E(X)$,
$\E(X^2)$, $\Var(X)$, and $\E(3X^2 - 2X + 1)$.

\begin{solution}
\[
  \E(X) = \int_0^4 x\cdot\frac{3}{32}x(4-x)\,\mathrm{d}x
  = \frac{3}{32}\int_0^4 (4x^2-x^3)\,\mathrm{d}x
  = \frac{3}{32}\left[\frac{4x^3}{3}-\frac{x^4}{4}\right]_0^4
  = \frac{3}{32}\!\left(\frac{256}{3} - 64\right)
  = \frac{3}{32}\cdot\frac{64}{3} = 2.
\]

\[
  \E(X^2) = \frac{3}{32}\int_0^4 (4x^3-x^4)\,\mathrm{d}x
  = \frac{3}{32}\left[x^4 - \frac{x^5}{5}\right]_0^4
  = \frac{3}{32}\!\left(256 - \frac{1024}{5}\right)
  = \frac{3}{32}\cdot\frac{256}{5} = \frac{24}{5} = 4.8.
\]

\[
  \Var(X) = \E(X^2) - [\E(X)]^2 = 4.8 - 4 = 0.8.
\]

By linearity of expectation:
\[
  \E(3X^2-2X+1) = 3\E(X^2) - 2\E(X) + 1 = 3(4.8) - 2(2) + 1 = 14.4 - 4 + 1 = 11.4.
\]
\end{solution}
\end{workedexample}

\subsection{Mode and Median of a Continuous Distribution}

\begin{definition}[Mode]
The \textbf{mode} of a continuous distribution is the value of $x$ at which
$f(x)$ is maximised.  It is found by solving $f'(x)=0$ and checking it is a
maximum.
\end{definition}

Note that a distribution may be \textbf{bimodal} (two local maxima) or even
\textbf{multimodal}.

\begin{workedexample}[--- Finding the Mode]
For $f(x) = \frac{3}{32}x(4-x)$, $0\le x\le 4$:
\begin{solution}
\[
  f'(x) = \frac{3}{32}(4-2x) = 0 \implies x = 2.
\]
Since $f''(2) = \frac{3}{32}(-2) < 0$, this is a maximum.  The mode is $x=2$.
\end{solution}
\end{workedexample}

\section{Finding the PDF and CDF of $Y = g(X)$}
\label{sec:transformation}

\subsection{Why Transformations Matter}

Often we know the distribution of $X$ but we want the distribution of a
function $Y = g(X)$.  For example:
\begin{itemize}
  \item If $X$ is the radius of a sphere (uniform), what is the distribution
        of its volume $Y = \frac{4}{3}\pi X^3$?
  \item If $X \sim \operatorname{Exp}(\lambda)$, what is the distribution of
        $Y = X^2$ (which arises in Weibull reliability models)?
\end{itemize}

\subsection{The CDF Method}

The most general and reliable approach is the \textbf{CDF method}:
\begin{enumerate}[label=\textbf{Step \arabic*:}]
  \item Write $F_Y(y) = \Prob(Y \le y) = \Prob(g(X) \le y)$.
  \item Rearrange the inequality to get bounds on $X$.
  \item Express in terms of $F_X$.
  \item Differentiate to get $f_Y(y) = F_Y'(y)$.
\end{enumerate}

\begin{workedexample}[--- Monotone Increasing Transformation]
Let $X$ have PDF $f_X(x) = 2x$, $0\le x\le 1$.  Find the PDF of $Y = X^3$.

\begin{solution}
\textbf{Step 1.}  The support of $Y$ is $[0,1]$.  For $y\in[0,1]$:
\[
  F_Y(y) = \Prob(Y\le y) = \Prob(X^3\le y) = \Prob\!\left(X\le y^{1/3}\right)
  = F_X\!\left(y^{1/3}\right) = \int_0^{y^{1/3}} 2t\,\mathrm{d}t = y^{2/3}.
\]

\textbf{Step 2.}  Differentiate:
\[
  f_Y(y) = \frac{\mathrm{d}}{\mathrm{d}y}\,y^{2/3} = \frac{2}{3}y^{-1/3},
  \quad 0 < y \le 1.
\]
\end{solution}
\end{workedexample}

\subsection{The Change-of-Variables Formula (Jacobian Method)}

For a \textbf{strictly monotone} transformation $y = g(x)$, there is a
slick formula:

\begin{theorem}[Change of Variables --- One Dimension]
Let $X$ have PDF $f_X$.  If $Y = g(X)$ where $g$ is strictly monotone and
differentiable, then
\[
  f_Y(y) = f_X\!\left(g^{-1}(y)\right)\left|\frac{\mathrm{d}}{\mathrm{d}y}g^{-1}(y)\right|.
\]
\end{theorem}

\begin{proof}
\textbf{Case: $g$ strictly increasing.}
\[
  F_Y(y) = \Prob(g(X)\le y) = \Prob\!\left(X\le g^{-1}(y)\right) = F_X\!\left(g^{-1}(y)\right).
\]
Differentiating: $f_Y(y) = f_X(g^{-1}(y))\cdot\frac{\mathrm{d}}{\mathrm{d}y}g^{-1}(y)$.
Since $g$ is increasing, $g^{-1}$ is increasing, so the derivative is
positive and equals the absolute value.

\textbf{Case: $g$ strictly decreasing.}
\[
  F_Y(y) = \Prob(g(X)\le y) = \Prob\!\left(X\ge g^{-1}(y)\right)
  = 1 - F_X\!\left(g^{-1}(y)\right).
\]
Differentiating: $f_Y(y) = -f_X(g^{-1}(y))\cdot\frac{\mathrm{d}}{\mathrm{d}y}g^{-1}(y)$.
Since $g$ is decreasing, $g^{-1}$ is decreasing, so $\frac{\mathrm{d}}{\mathrm{d}y}g^{-1}(y)<0$
and the product is positive.  Taking the absolute value unifies both cases.
\end{proof}

\begin{workedexample}[--- Applying the Jacobian Formula]
Let $X \sim U(0,1)$.  Find the PDF of $Y = -\ln X$.

\begin{solution}
$g(x) = -\ln x$ is strictly decreasing on $(0,1)$.
The inverse: $x = g^{-1}(y) = e^{-y}$.
The support of $Y$: as $x\in(0,1)$, $-\ln x \in (0,\infty)$, so $y\in(0,\infty)$.

\[
  f_Y(y) = f_X(e^{-y})\cdot\left|\frac{\mathrm{d}}{\mathrm{d}y}e^{-y}\right|
  = 1 \cdot e^{-y} = e^{-y}, \quad y > 0.
\]

We recognise $f_Y(y) = e^{-y}$ as the $\operatorname{Exp}(1)$ distribution.
This is a famous result: if $U\sim U(0,1)$ then $-\ln U \sim \operatorname{Exp}(1)$.
It is used in \textbf{inverse transform sampling} to simulate exponential random
variables from uniform ones.
\end{solution}
\end{workedexample}

\begin{realworld}[: Simulation]
The inverse transform method allows computers to generate samples from any
continuous distribution using only uniform random numbers (which computers
can generate efficiently).  If $U\sim U(0,1)$, then $X = F^{-1}(U)$ has CDF
$F$.  This is foundational in Monte Carlo simulation, used in finance,
particle physics, and climate modelling.
\end{realworld}

\begin{workedexample}[--- Non-monotone Transformation]
Let $X\sim U(-1,1)$.  Find the PDF of $Y = X^2$.

\begin{solution}
The function $g(x)=x^2$ is \emph{not} monotone on $(-1,1)$.  We use the CDF method.

Support of $Y$: $y\in[0,1]$.  For $y\in(0,1)$:
\begin{align*}
  F_Y(y) &= \Prob(X^2 \le y) = \Prob(-\sqrt{y}\le X\le\sqrt{y})
  = \int_{-\sqrt{y}}^{\sqrt{y}} \frac{1}{2}\,\mathrm{d}x = \sqrt{y}.
\end{align*}
Differentiating:
\[
  f_Y(y) = \frac{1}{2\sqrt{y}}, \quad 0 < y < 1.
\]
This is the $\operatorname{Beta}(\tfrac{1}{2},1)$ distribution (a specific case of the
Beta family you will encounter at university).
\end{solution}
\end{workedexample}

\section{Advanced Exploration: The Normal Distribution Revisited}

\begin{exploration}
\textbf{Why is the Normal Distribution Normal?}

The normal distribution $X\sim N(\mu,\sigma^2)$ has PDF
\[
  f(x) = \frac{1}{\sigma\sqrt{2\pi}}\exp\!\left(-\frac{(x-\mu)^2}{2\sigma^2}\right).
\]
\begin{enumerate}
  \item Verify that $\int_{-\infty}^{\infty} e^{-x^2/2}\,\mathrm{d}x = \sqrt{2\pi}$
        using the classic trick $\left(\int e^{-x^2/2}\mathrm{d}x\right)^2 = \int\!\int e^{-(x^2+y^2)/2}\mathrm{d}x\,\mathrm{d}y$
        and polar coordinates.  This integral is one of the most beautiful in all mathematics.
  \item Show that if $X\sim N(0,1)$, then $Y=X^2\sim\chi^2(1)$ using the transformation formula.
        (This is the seed of the chi-squared distribution you will study in Chapter 10.)
  \item The \textbf{Central Limit Theorem} says that the standardised sum of iid random
        variables converges to $N(0,1)$.  Research what ``converges in distribution'' means
        and why the normal distribution is the unique stable limit.
\end{enumerate}
\end{exploration}

\section{Common Mistakes in Chapter 8}

\begin{warningbox}
\begin{enumerate}[label=\textbf{\arabic*.}]
  \item \textbf{Confusing $f(x)$ with probability.}  $f(x)$ can exceed 1;
        only integrals give probabilities.
  \item \textbf{Forgetting limits of integration.}  Always identify the
        support of $f$ and integrate only over that region.
  \item \textbf{Using the Jacobian formula for non-monotone functions.}
        The formula only applies when $g$ is strictly monotone.  For others,
        use the CDF method.
  \item \textbf{Forgetting the absolute value} in the Jacobian formula.
        The PDF must be non-negative.
  \item \textbf{Forgetting to state the support} of the transformed variable.
        Always specify where $f_Y(y) > 0$.
\end{enumerate}
\end{warningbox}

\section{Summary of Chapter 8}

\begin{summarybox}
\begin{description}[style=nextline]
  \item[PDF] $f(x)\ge 0$, $\int_{-\infty}^\infty f(x)\,\mathrm{d}x=1$.
    Probabilities are areas: $\Prob(a\le X\le b)=\int_a^b f(x)\,\mathrm{d}x$.
  \item[CDF] $F(x)=\int_{-\infty}^x f(t)\,\mathrm{d}t$.
    Key relationships: $f=F'$, $\Prob(a\le X\le b)=F(b)-F(a)$.
  \item[Expectation] $\E(g(X))=\int_{-\infty}^\infty g(x)f(x)\,\mathrm{d}x$ (LOTUS).
  \item[Variance] $\Var(X)=\E(X^2)-[\E(X)]^2$.
  \item[Transformations] CDF method (general) or Jacobian formula (monotone only).
  \item[Key property] $\Prob(X=a)=0$ for all $a$.
  \item[Mode] Value maximising $f(x)$.
  \item[Median $m$] $F(m)=0.5$.
\end{description}
\end{summarybox}

% ══════════════════════════════════════════════════════════
%  CHAPTER 9: INFERENTIAL STATISTICS
% ══════════════════════════════════════════════════════════
\chapter{Inferential Statistics}
\label{ch:infer}

\section{Introduction}

Statistical inference is the art and science of drawing conclusions about an
\emph{unknown population} from a \emph{limited sample}.  You have already
seen inference using the normal distribution (known variance) and the Poisson
distribution.  This chapter introduces two new settings of profound practical
importance:

\begin{enumerate}
  \item Situations where the \textbf{population variance is unknown} and
        the sample is small---we must use the $t$-distribution.
  \item Situations where we compare \textbf{two populations} (are their
        means equal?) or analyse \textbf{paired data} (before/after measurements).
\end{enumerate}

These tools are among the most widely used in medicine, psychology,
economics, and engineering.  Every clinical trial that concludes a new drug
is effective, every study that finds a wage gap between groups, every
experiment that detects a new elementary particle---all rely on the machinery
of this chapter.

\section{The $t$-Distribution}
\label{sec:tdist}

\subsection{The Problem with Unknown Variance}

Suppose $X_1,X_2,\ldots,X_n \iid N(\mu,\sigma^2)$ and we want to make
inferences about $\mu$.  If $\sigma^2$ is \emph{known}, the pivot
\[
  Z = \frac{\bar{X} - \mu}{\sigma/\sqrt{n}} \sim N(0,1)
\]
gives exact inference.

But in practice, $\sigma^2$ is \emph{almost never known}.  We estimate it
by the sample variance
\[
  S^2 = \frac{1}{n-1}\sum_{i=1}^n (X_i - \bar{X})^2.
\]
The natural thing is to replace $\sigma$ by $S$, giving the \textbf{$t$-statistic}:
\[
  T = \frac{\bar{X}-\mu}{S/\sqrt{n}}.
\]

What is the distribution of $T$?  It is \emph{not} normal, because $S$ is
itself random and introduces extra variability.

\subsection{Student's $t$-Distribution: History and Definition}

\begin{historicalbox}
William Sealy Gosset (1876--1937) was a chemist working for Guinness Brewery
in Dublin.  He derived the distribution of $T$ in 1908, but Guinness forbade
employees from publishing under their own names (to protect trade secrets).
He published under the pseudonym ``Student''---hence \emph{Student's
$t$-distribution}.  His work solved the fundamental problem of inference with
small samples and unknown variance.  Ronald A.\ Fisher later placed it in a
broader theoretical framework.
\end{historicalbox}

\begin{definition}[Student's $t$-Distribution]
If $X_1,\ldots,X_n \iid N(\mu,\sigma^2)$, then
\[
  T = \frac{\bar{X}-\mu}{S/\sqrt{n}} \sim t(\nu),
\]
where $\nu = n-1$ is the number of \textbf{degrees of freedom} and
$S^2 = \frac{1}{n-1}\sum(X_i-\bar{X})^2$ is the sample variance.
\end{definition}

\begin{intuition}
\emph{Degrees of freedom} measure how much ``independent information'' is
available to estimate variability.  Given $n$ observations, once we have
fixed $\bar{X}$, only $n-1$ of the deviations $X_i-\bar{X}$ are free to
vary (the last one is determined by the constraint
$\sum(X_i-\bar{X})=0$).  Hence $\nu = n-1$.
\end{intuition}

\subsection{Shape of the $t$-Distribution}

The $t(\nu)$ distribution:
\begin{itemize}
  \item Is \textbf{symmetric and bell-shaped}, centred at 0.
  \item Has \textbf{heavier tails} than the normal---reflects the extra
        uncertainty from estimating $\sigma$.
  \item Converges to $N(0,1)$ as $\nu\to\infty$: with large samples,
        $S\approx\sigma$ and the extra uncertainty disappears.
  \item For $\nu=1$, it is the \textbf{Cauchy distribution} (no finite mean!).
\end{itemize}

\begin{figure}[h]
\centering
\begin{tikzpicture}
\begin{axis}[
  width=13cm, height=6cm,
  xlabel={$t$}, ylabel={Density},
  xmin=-5, xmax=5, ymin=0, ymax=0.42,
  axis lines=left,
  legend pos=north east,
  legend style={font=\small},
  title={Student's $t$-distributions vs the Standard Normal},
  title style={font=\bfseries},
  samples=200,
]
% Normal
\addplot[thick, PrimaryBlue, domain=-5:5]
  {exp(-x^2/2)/sqrt(2*pi)};
\addlegendentry{$N(0,1)$}
% t(1) = Cauchy
\addplot[thick, AccentRed, domain=-5:5]
  {1/(pi*(1+x^2))};
\addlegendentry{$t(1)$}
% t(3)
\addplot[thick, AccentOrange, domain=-5:5]
  {(1 + x^2/3)^(-2) * (8/(3*pi*sqrt(3)))};
\addlegendentry{$t(3)$}
% t(10)
\addplot[thick, AccentTeal, domain=-5:5]
  {(1 + x^2/10)^(-5.5) * 3.4914/10^(0.5)};
\addlegendentry{$t(10)$}
\end{axis}
\end{tikzpicture}
\caption{The $t$-distribution has heavier tails than the normal.  As
  degrees of freedom increase, it approaches the standard normal.}
\label{fig:tdist}
\end{figure}

\subsection{Using $t$-Tables}

Since $T\sim t(\nu)$ is symmetric about 0, tables typically give
\emph{one-tailed} critical values.  The value $t_{\alpha,\nu}$ satisfies
$\Prob(T > t_{\alpha,\nu}) = \alpha$.

For a \textbf{two-tailed} test at significance level $\alpha$, we reject
$H_0$ when $|T| > t_{\alpha/2,\nu}$.

\begin{examtrap}
Always look up the correct degrees of freedom.  For a single sample of
size $n$, use $\nu = n-1$.  For two independent samples of sizes $n_1$
and $n_2$, use $\nu = n_1+n_2-2$ (if variances are assumed equal).
\end{examtrap}

\section{Hypothesis Tests for the Difference in Means}
\label{sec:twosample}

\subsection{Setting}

Suppose we have two independent populations:
\begin{itemize}
  \item Population 1: $X_{11},\ldots,X_{1n_1}\iid N(\mu_1,\sigma^2)$, with
        sample mean $\bar{X}_1$ and sample variance $S_1^2$.
  \item Population 2: $X_{21},\ldots,X_{2n_2}\iid N(\mu_2,\sigma^2)$, with
        sample mean $\bar{X}_2$ and sample variance $S_2^2$.
\end{itemize}

We assume \textbf{equal but unknown variance} $\sigma^2$ (this is a
testable assumption---we will discuss this later).

\subsection{The Pooled Variance Estimator}

We combine (pool) the two sample variances into a single estimate of $\sigma^2$:
\[
  S_p^2 = \frac{(n_1-1)S_1^2 + (n_2-1)S_2^2}{n_1+n_2-2}.
\]

\begin{intuition}
$S_p^2$ is a \emph{weighted average} of the two sample variances, with
weights proportional to the degrees of freedom.  If one sample is much
larger than the other, its variance estimate gets more weight.
\end{intuition}

\subsection{The Two-Sample $t$-Test}

\begin{proposition}[Two-Sample $t$-Statistic]
Under $H_0: \mu_1 = \mu_2$ (equivalently $\mu_1-\mu_2 = 0$),
\[
  T = \frac{\bar{X}_1 - \bar{X}_2}{S_p\sqrt{\dfrac{1}{n_1}+\dfrac{1}{n_2}}}
  \sim t(n_1+n_2-2).
\]
\end{proposition}

\begin{proof}[Sketch]
Since $\bar{X}_1\sim N(\mu_1,\sigma^2/n_1)$ and $\bar{X}_2\sim N(\mu_2,\sigma^2/n_2)$
independently,
\[
  \bar{X}_1 - \bar{X}_2 \sim N\!\left(\mu_1-\mu_2,\;\sigma^2\!\left(\frac{1}{n_1}+\frac{1}{n_2}\right)\right).
\]
Under $H_0$, the numerator is $N(0,\sigma^2(1/n_1+1/n_2))$.  The
denominator $S_p\sqrt{1/n_1+1/n_2}$ estimates the standard deviation of the
numerator.  One can show $(n_1+n_2-2)S_p^2/\sigma^2\sim\chi^2(n_1+n_2-2)$
independently of the numerator, giving the $t$ distribution by definition.
\end{proof}

\subsection{Hypothesis Testing Procedure}

\begin{enumerate}[label=\textbf{Step \arabic*:}]
  \item \textbf{State hypotheses:} $H_0:\mu_1=\mu_2$ vs
        $H_1:\mu_1\ne\mu_2$ (two-tailed), or
        $H_1:\mu_1>\mu_2$ / $H_1:\mu_1<\mu_2$ (one-tailed).
  \item \textbf{State assumptions:} Both populations normal, equal
        (unknown) variances, independent samples.
  \item \textbf{Compute the pooled variance:}
        $S_p^2 = \frac{(n_1-1)S_1^2+(n_2-1)S_2^2}{n_1+n_2-2}$.
  \item \textbf{Compute the test statistic:}
        $t_{\text{calc}} = \frac{\bar{x}_1-\bar{x}_2}{s_p\sqrt{1/n_1+1/n_2}}$.
  \item \textbf{Find critical value:} $t_{\alpha/2,\,n_1+n_2-2}$ (two-tailed)
        from tables.
  \item \textbf{Compare and conclude:} Reject $H_0$ if $|t_{\text{calc}}|>t_{\text{crit}}$.
  \item \textbf{Interpret} in context.
\end{enumerate}

\begin{workedexample}[--- Two-Sample $t$-Test]
A researcher wants to know whether a new fertiliser increases crop yield.
Independent samples give:
\begin{center}
\begin{tabular}{lccc}
\toprule
  & $n$ & $\bar{x}$ & $s^2$ \\
\midrule
Control   & 10 & 45.2 & 12.4 \\
Treatment & 12 & 49.8 & 14.6 \\
\bottomrule
\end{tabular}
\end{center}
Test at the 5\% significance level whether the treatment increases yield.

\begin{solution}
\textbf{Hypotheses:} $H_0:\mu_C=\mu_T$ vs $H_1:\mu_T>\mu_C$ (one-tailed).

\textbf{Pooled variance:}
\[
  s_p^2 = \frac{9(12.4)+11(14.6)}{10+12-2}
  = \frac{111.6+160.6}{20}
  = \frac{272.2}{20} = 13.61.
\]
So $s_p = \sqrt{13.61} = 3.690$.

\textbf{Test statistic:}
\[
  t = \frac{49.8-45.2}{3.690\sqrt{1/12+1/10}}
  = \frac{4.6}{3.690\sqrt{0.1833}}
  = \frac{4.6}{3.690\times 0.4282}
  = \frac{4.6}{1.580} = 2.911.
\]

\textbf{Critical value:} $\nu=20$, one-tailed at 5\%: $t_{0.05,20}=1.725$.

\textbf{Decision:} $2.911 > 1.725$, so we reject $H_0$.

\textbf{Conclusion:} There is significant evidence at the 5\% level that
the fertiliser increases yield.
\end{solution}
\end{workedexample}

\section{Paired $t$-Tests}
\label{sec:paired}

\subsection{When to Use Paired Tests}

The two-sample test assumes the two samples are \emph{independent}.  But
sometimes the data are \textbf{naturally paired}:
\begin{itemize}
  \item Blood pressure of the same patients before and after treatment.
  \item Test scores of students taught by two methods, where each student
        takes both.
  \item Measurements on left and right limbs.
\end{itemize}

In these cases, observations within a pair are \emph{correlated}, and
treating them as independent loses information.  The paired $t$-test exploits
the pairing by reducing the problem to a \emph{one-sample} test on the
differences.

\subsection{The Paired $t$-Statistic}

Let $(X_{1i}, X_{2i})$ be the $i$-th pair.  Define the difference:
\[
  D_i = X_{1i} - X_{2i}, \quad i=1,\ldots,n.
\]

If the underlying distributions are normal, then $D_i\sim N(\mu_D,\sigma_D^2)$
for some unknown $\mu_D$ and $\sigma_D^2$.  The test reduces to a one-sample
$t$-test on $D_1,\ldots,D_n$:
\[
  H_0: \mu_D = 0 \quad \text{vs} \quad H_1: \mu_D \ne 0.
\]
The test statistic is
\[
  T = \frac{\bar{D}}{S_D/\sqrt{n}} \sim t(n-1) \quad \text{under } H_0,
\]
where $\bar{D} = \frac{1}{n}\sum D_i$ and
$S_D^2 = \frac{1}{n-1}\sum(D_i-\bar{D})^2$.

\begin{intuition}
By looking only at the differences, we \emph{eliminate the variability
between subjects}.  If patient A naturally has higher blood pressure than
patient B, this fact doesn't affect the difference caused by the drug.
Pairing removes this ``between-subject'' noise, making the test more
powerful.
\end{intuition}

\begin{workedexample}[--- Paired $t$-Test]
Eight athletes had their sprint times (in seconds) measured before and after
a training programme:
\begin{center}
\begin{tabular}{lrrrrrrrr}
\toprule
Athlete & 1 & 2 & 3 & 4 & 5 & 6 & 7 & 8 \\
\midrule
Before  & 12.4 & 11.8 & 13.2 & 12.9 & 11.5 & 12.7 & 13.0 & 12.1 \\
After   & 12.0 & 11.4 & 12.9 & 12.5 & 11.3 & 12.2 & 12.4 & 11.8 \\
$D_i$   & 0.4  & 0.4  & 0.3  & 0.4  & 0.2  & 0.5  & 0.6  & 0.3  \\
\bottomrule
\end{tabular}
\end{center}
Test at the 1\% level whether the programme reduces sprint times.

\begin{solution}
\textbf{Hypotheses:} $H_0:\mu_D=0$ vs $H_1:\mu_D>0$ (one-tailed, since a
positive difference means before $>$ after, i.e.\ times decreased).

\[
  \bar{d} = \frac{0.4+0.4+0.3+0.4+0.2+0.5+0.6+0.3}{8} = \frac{3.1}{8} = 0.3875.
\]
\[
  s_D^2 = \frac{1}{7}\left[\sum d_i^2 - 8\bar{d}^2\right]
  = \frac{1}{7}\left[(0.16+0.16+0.09+0.16+0.04+0.25+0.36+0.09) - 8(0.3875)^2\right].
\]
$\sum d_i^2 = 1.31$.  $8(0.3875)^2 = 8(0.1502) = 1.2012$.
\[
  s_D^2 = \frac{1.31-1.2012}{7} = \frac{0.1088}{7} = 0.01554,\quad s_D=0.1247.
\]
\[
  t = \frac{0.3875}{0.1247/\sqrt{8}} = \frac{0.3875}{0.04410} = 8.786.
\]

Critical value: $t_{0.01,7}=2.998$.

Since $8.786 > 2.998$, we reject $H_0$.  There is very strong evidence that
the programme reduces sprint times.
\end{solution}
\end{workedexample}

\begin{warningbox}
Never use a two-sample $t$-test when data are paired.  Doing so ignores the
correlation between paired observations, inflates the variance estimate, and
reduces the power of the test---you may fail to detect a real effect.
Similarly, never use a paired $t$-test when samples are truly independent.
\end{warningbox}

\section{Confidence Intervals for the Mean (Small Sample)}
\label{sec:ci-mean}

\subsection{Construction}

A \textbf{$(1-\alpha)\times 100\%$ confidence interval} for the population
mean $\mu$, when $\sigma^2$ is unknown and the population is normal, is:
\[
  \bar{x} \pm t_{\alpha/2,\,n-1} \cdot \frac{s}{\sqrt{n}}.
\]

\begin{intuition}
The CI captures the \emph{plausible range} of values for $\mu$ given the
data.  The factor $t_{\alpha/2,n-1}$ is wider than the normal $z_{\alpha/2}$
because we have estimated $\sigma$: we must account for that extra uncertainty.
As $n\to\infty$, $t\to z$ and the CI based on the normal approximation
becomes valid.
\end{intuition}

\begin{keyresult}
\textbf{Correct interpretation:}
``We are 95\% confident that the true population mean lies between $L$ and $U$.''

This means: if we repeated the sampling process many times and computed this
CI each time, 95\% of the resulting intervals would contain the true $\mu$.
The interval is random; the parameter is fixed.  It is \emph{wrong} to say
``there is a 95\% probability that $\mu$ is in $(L,U)$'' once the data are observed.
\end{keyresult}

\begin{figure}[h]
\centering
\begin{tikzpicture}[scale=0.9]
% Axes
\draw[->] (-0.5,0) -- (10.5,0) node[right] {$\mu$};
\node[below] at (5,-0.3) {True $\mu$};
\draw[dashed,MidGrey] (5,-0.2) -- (5,8.5);

% 20 simulated CIs
\pgfmathsetseed{42}
\foreach \i in {1,...,20} {
  \pgfmathsetmacro{\center}{5 + 0.7*rand}
  \pgfmathsetmacro{\half}{0.8 + 0.3*rnd}
  \pgfmathsetmacro{\left}{\center - \half}
  \pgfmathsetmacro{\right}{\center + \half}
  \pgfmathsetmacro{\ypos}{0.38*\i}
  % Colour red if interval misses mu=5
  \ifdim \left pt < 5pt
    \ifdim \right pt > 5pt
      \draw[PrimaryBlue, thick] (\left,\ypos) -- (\right,\ypos);
      \fill[PrimaryBlue] (\center,\ypos) circle (1.5pt);
    \else
      \draw[AccentRed, thick] (\left,\ypos) -- (\right,\ypos);
      \fill[AccentRed] (\center,\ypos) circle (1.5pt);
    \fi
  \else
    \draw[AccentRed, thick] (\left,\ypos) -- (\right,\ypos);
    \fill[AccentRed] (\center,\ypos) circle (1.5pt);
  \fi
}
\node[right, font=\small] at (10.5,4) {\color{PrimaryBlue}Contains $\mu$};
\node[right, font=\small] at (10.5,3) {\color{AccentRed}Misses $\mu$};
\end{tikzpicture}
\caption{Simulated 95\% confidence intervals across 20 repeated samples.
  Approximately 95\% (blue) capture the true $\mu$; a few (red) do not.
  The parameter is fixed; the intervals are random.}
\label{fig:ci-sim}
\end{figure}

\begin{workedexample}[--- Constructing a CI]
A sample of 15 measurements of a chemical concentration gives
$\bar{x}=8.32$ and $s=1.45$.  Construct a 99\% CI for the true mean
concentration.

\begin{solution}
$n=15$, $\nu=14$, $\alpha=0.01$, $t_{0.005,14}=2.977$.
\[
  \bar{x}\pm t_{0.005,14}\cdot\frac{s}{\sqrt{n}}
  = 8.32\pm 2.977\times\frac{1.45}{\sqrt{15}}
  = 8.32\pm 2.977\times 0.3744
  = 8.32\pm 1.114.
\]
The 99\% CI is $(7.206,\;9.434)$.

\medskip
\textit{Interpretation:} We are 99\% confident that the true mean
concentration lies between 7.206 and 9.434 units.
\end{solution}
\end{workedexample}

\section{Confidence Intervals for the Difference in Means}
\label{sec:ci-diff}

\subsection{Independent Samples (Pooled)}

A $(1-\alpha)\times 100\%$ CI for $\mu_1-\mu_2$ (with equal unknown
variances) is:
\[
  (\bar{x}_1-\bar{x}_2) \pm t_{\alpha/2,\,n_1+n_2-2}\cdot s_p\sqrt{\frac{1}{n_1}+\frac{1}{n_2}}.
\]

\subsection{Paired Data}

A $(1-\alpha)\times 100\%$ CI for $\mu_D = \mu_1 - \mu_2$ is:
\[
  \bar{d} \pm t_{\alpha/2,\,n-1}\cdot\frac{s_D}{\sqrt{n}}.
\]

\begin{workedexample}[--- CI for Difference in Means]
Using the data from the fertiliser example (Section~\ref{sec:twosample}):
$n_1=10$, $\bar{x}_1=45.2$, $n_2=12$, $\bar{x}_2=49.8$, $s_p^2=13.61$.
Construct a 95\% CI for $\mu_T-\mu_C$.

\begin{solution}
$\nu=20$, $t_{0.025,20}=2.086$.
\[
  (\bar{x}_2-\bar{x}_1)\pm 2.086\times s_p\sqrt{\frac{1}{n_1}+\frac{1}{n_2}}
  = 4.6\pm 2.086\times 3.690\times 0.4282
  = 4.6\pm 3.296.
\]
95\% CI for $\mu_T-\mu_C$: $(1.304,\;7.896)$.

Since the interval does not contain 0, there is significant evidence of a
positive treatment effect (consistent with the one-sided test).
\end{solution}
\end{workedexample}

\subsection{Type I and Type II Errors}

\begin{definition}[Type I and Type II Errors]
\begin{center}
\begin{tabular}{lcc}
\toprule
& $H_0$ is true & $H_0$ is false \\
\midrule
Reject $H_0$    & \textbf{Type I error} ($\alpha$) & Correct (power $=1-\beta$) \\
Do not reject   & Correct & \textbf{Type II error} ($\beta$) \\
\bottomrule
\end{tabular}
\end{center}
\end{definition}

\begin{itemize}
  \item The \textbf{significance level} $\alpha = \Prob(\text{Type I error})$ is set
        by the researcher before collecting data.
  \item The \textbf{power} $1-\beta = \Prob(\text{reject }H_0\mid H_0\text{ false})$
        measures the test's sensitivity to a real effect.
  \item There is a fundamental trade-off: reducing $\alpha$ (being more
        conservative) increases $\beta$ (missing real effects).
\end{itemize}

\begin{realworld}[: Clinical Trials]
In drug approval, a Type I error means approving an ineffective drug
(with possible side effects and wasted resources).  A Type II error means
rejecting a drug that actually works (denying patients a beneficial treatment).
Regulators typically set $\alpha=0.05$ as a balance, but for life-threatening
diseases, higher power (lower $\beta$) may be prioritised.
\end{realworld}

\section{Advanced Exploration: Welch's $t$-Test and Normality}

\begin{exploration}
\textbf{What if variances are not equal?}

The pooled $t$-test assumes equal population variances.  If this is suspect:
\begin{enumerate}
  \item \textbf{Welch's $t$-test} uses a modified test statistic with
        approximate degrees of freedom given by the Welch--Satterthwaite
        equation:
        \[
          \nu \approx \frac{(s_1^2/n_1 + s_2^2/n_2)^2}
                            {\frac{(s_1^2/n_1)^2}{n_1-1}+\frac{(s_2^2/n_2)^2}{n_2-1}}.
        \]
        This is the default in most statistical software.  Research why
        this formula arises.
  \item \textbf{Levene's test} or \textbf{Bartlett's test} formally tests
        $H_0:\sigma_1^2=\sigma_2^2$.  But many statisticians argue you should
        always use Welch's test---it is robust to both equal and unequal variances.
  \item The Central Limit Theorem means that for large $n$, the $t$-test is
        approximately valid even for non-normal populations.  Research what
        ``sufficiently large'' means and when normality matters.
\end{enumerate}
\end{exploration}

\section{Summary of Chapter 9}

\begin{summarybox}
\begin{description}[style=nextline]
  \item[$t$-distribution] $T=\frac{\bar{X}-\mu}{S/\sqrt{n}}\sim t(n-1)$.
    Symmetric, heavier tails than normal; approaches $N(0,1)$ as $n\to\infty$.
  \item[Two-sample $t$-test] Pooled variance $s_p^2$; $\nu=n_1+n_2-2$.
    Use when samples are independent and variances assumed equal.
  \item[Paired $t$-test] Compute differences $D_i=X_{1i}-X_{2i}$;
    one-sample $t$-test on $D_i$; $\nu=n-1$.  Use when data are naturally paired.
  \item[Confidence interval] $\bar{x}\pm t_{\alpha/2,\nu}\cdot\frac{s}{\sqrt{n}}$ (one sample);
    analogous formulas for differences.
  \item[Interpretation] CI captures plausible range for the parameter.
    Significance test controls Type I error probability $\alpha$.
\end{description}
\end{summarybox}

% ══════════════════════════════════════════════════════════
%  CHAPTER 10: CHI-SQUARED TESTS
% ══════════════════════════════════════════════════════════
\chapter{Chi-Squared Tests}
\label{ch:chi}

\section{Introduction}

The $t$-test asks: is the population mean a particular value, or are two
means equal?  But often we want to compare \emph{entire distributions}, or
test whether two \emph{categorical} variables are associated.  For these
questions, we need the chi-squared ($\chi^2$) test.

Chi-squared tests arise in:
\begin{itemize}
  \item \textbf{Genetics:} Do observed genotype frequencies match Mendelian
        ratios?
  \item \textbf{Medicine:} Is there an association between smoking and lung cancer?
  \item \textbf{Marketing:} Does purchasing behaviour depend on age group?
  \item \textbf{Quality control:} Does the defect rate follow an assumed distribution?
\end{itemize}

\section{The Chi-Squared Distribution}
\label{sec:chi-dist}

\begin{definition}[Chi-Squared Distribution]
If $Z_1,Z_2,\ldots,Z_\nu \iid N(0,1)$, then
\[
  Q = Z_1^2 + Z_2^2 + \cdots + Z_\nu^2 \sim \chi^2(\nu).
\]
The parameter $\nu$ is the degrees of freedom.
\end{definition}

Properties:
\begin{itemize}
  \item Support: $(0,\infty)$ (it is a sum of squares---always non-negative).
  \item $\E(\chi^2(\nu)) = \nu$.
  \item $\Var(\chi^2(\nu)) = 2\nu$.
  \item Shape: skewed right for small $\nu$; approaches normal as $\nu\to\infty$.
\end{itemize}

\begin{figure}[h]
\centering
\begin{tikzpicture}
\begin{axis}[
  width=13cm, height=6cm,
  xlabel={$x$}, ylabel={Density},
  xmin=0, xmax=20, ymin=0, ymax=0.38,
  axis lines=left,
  legend pos=north east,
  legend style={font=\small},
  title={Chi-squared distributions for different degrees of freedom},
  title style={font=\bfseries},
  samples=200,
  restrict y to domain=0:0.5,
]
\addplot[thick, PrimaryBlue, domain=0.1:20]
  {(x^(1/2-1)*exp(-x/2))/(2^(1/2)*0.8862)};
\addlegendentry{$\chi^2(1)$}
\addplot[thick, AccentOrange, domain=0.1:20]
  {x^(2-1)*exp(-x/2)/(2^2*1)};
\addlegendentry{$\chi^2(4)$}
\addplot[thick, AccentTeal, domain=0.1:20]
  {x^(5-1)*exp(-x/2)/(2^5*24)};
\addlegendentry{$\chi^2(10)$}
\addplot[thick, AccentRed, domain=0.1:20]
  {x^(8-1)*exp(-x/2)/(2^8*5040)};
\addlegendentry{$\chi^2(16)$ (approx)}
\end{axis}
\end{tikzpicture}
\caption{The $\chi^2$ distribution is right-skewed for small $\nu$ and
  becomes more symmetric as $\nu$ increases.}
\label{fig:chisq}
\end{figure}

\section{The Chi-Squared Test Statistic}

The general form of the \textbf{Pearson chi-squared statistic} is:
\[
  X^2 = \sum_{\text{all cells}} \frac{(O-E)^2}{E},
\]
where $O$ is the \textbf{observed frequency} and $E$ is the
\textbf{expected frequency} under $H_0$.

\begin{intuition}
Each term $\frac{(O-E)^2}{E}$ measures how much a single cell deviates from
expectation, \emph{relative to} the expected count.  A deviation of 5 when
$E=100$ is small; the same deviation when $E=5$ is enormous.  Squaring makes
all terms positive and penalises large deviations.  If $H_0$ is true,
$X^2\approx\chi^2(\nu)$ for appropriate $\nu$.
\end{intuition}

\begin{warningbox}
The chi-squared approximation requires \textbf{expected frequencies of at
least 5} in each cell.  If some expected frequencies are smaller, merge
adjacent categories until all $E\ge 5$.  Small expected frequencies make
the $\chi^2$ approximation unreliable.
\end{warningbox}

\section{Goodness of Fit for Discrete Distributions}
\label{sec:gof-discrete}

\subsection{Setting and Hypotheses}

We have $n$ observations falling into $k$ categories.  We hypothesise that
the probability of each category is $p_1,p_2,\ldots,p_k$ (with $\sum p_i=1$).

\[
  H_0: \text{data follow the specified distribution}
  \quad\text{vs}\quad
  H_1: \text{they do not.}
\]

Expected frequencies: $E_i = np_i$.

The test statistic $X^2 = \sum_{i=1}^k \frac{(O_i-E_i)^2}{E_i}$ follows
$\chi^2(k-1-m)$ approximately under $H_0$, where $m$ is the number of
parameters estimated from the data.

\begin{itemize}
  \item No parameters estimated: $\nu = k-1$.
  \item One parameter estimated (e.g.\ $\lambda$ for Poisson): $\nu = k-2$.
\end{itemize}

\begin{workedexample}[--- Testing a Discrete Uniform Distribution]
A die is rolled 120 times.  Observed frequencies:
\begin{center}
\begin{tabular}{lcccccc}
\toprule
Face      & 1  & 2  & 3  & 4  & 5  & 6 \\
\midrule
Observed  & 24 & 18 & 22 & 16 & 21 & 19 \\
\bottomrule
\end{tabular}
\end{center}
Test at 5\% whether the die is fair.

\begin{solution}
$H_0$: each face equally likely ($p_i=1/6$).

Expected frequencies: $E_i = 120/6 = 20$ for all $i$.

\[
  X^2 = \frac{(24-20)^2}{20}+\frac{(18-20)^2}{20}+\frac{(22-20)^2}{20}
       +\frac{(16-20)^2}{20}+\frac{(21-20)^2}{20}+\frac{(19-20)^2}{20}
\]
\[
  = \frac{16+4+4+16+1+1}{20} = \frac{42}{20} = 2.1.
\]

Degrees of freedom: $\nu = 6-1 = 5$.  Critical value: $\chi^2_{0.05,5}=11.07$.

Since $2.1 < 11.07$, we do not reject $H_0$.  There is no significant
evidence that the die is unfair.
\end{solution}
\end{workedexample}

\begin{workedexample}[--- Fitting a Poisson Distribution]
The number of defects per item in a sample of 200 items:
\begin{center}
\begin{tabular}{lrrrrr}
\toprule
Defects $x$ & 0 & 1 & 2 & 3 & $\ge 4$ \\
\midrule
Observed    & 83 & 72 & 31 & 10 & 4 \\
\bottomrule
\end{tabular}
\end{center}
Test whether defects follow a Poisson distribution.

\begin{solution}
\textbf{Estimate $\lambda$:}
$\hat{\lambda} = \bar{x} = (0\times 83+1\times 72+2\times 31+3\times 10+4\times 4)/200$
$= (0+72+62+30+16)/200 = 180/200 = 0.9$.

\textbf{Expected frequencies} with $\hat{\lambda}=0.9$, $n=200$:
\begin{center}
\begin{tabular}{lrrrr}
\toprule
$x$ & $P(X=x)$ & $E_x$ & $O_x$ & $(O-E)^2/E$ \\
\midrule
0      & $e^{-0.9}=0.4066$         & 81.31 & 83 & $\frac{(1.69)^2}{81.31}=0.035$ \\
1      & $0.9e^{-0.9}=0.3659$      & 73.18 & 72 & $0.019$ \\
2      & $\frac{0.81}{2}e^{-0.9}=0.1647$ & 32.93 & 31 & $0.113$ \\
3      & $\frac{0.729}{6}e^{-0.9}=0.0494$ & 9.88  & 10 & $0.001$ \\
$\ge4$ & $1-\text{above}=0.0134$   & 2.70  & 4  & merged below \\
\bottomrule
\end{tabular}
\end{center}
Merge $x\ge 3$ (since $E_{\ge 4}=2.70 < 5$):
$O_{\ge 3}=14$, $E_{\ge 3}=9.88+2.70=12.58$.
$(14-12.58)^2/12.58 = 0.160$.

\[
  X^2 = 0.035 + 0.019 + 0.113 + 0.160 = 0.327.
\]
Degrees of freedom: $k=4$ categories, 1 parameter estimated, so $\nu=4-1-1=2$.

Critical value: $\chi^2_{0.05,2}=5.991$.  Since $0.327<5.991$, we do not reject $H_0$.
A Poisson distribution with $\lambda=0.9$ fits the data well.
\end{solution}
\end{workedexample}

\section{Goodness of Fit for Continuous Distributions}
\label{sec:gof-cont}

For continuous data, we \emph{group} observations into class intervals and
apply the same chi-squared test.  Expected frequencies are computed using
the assumed CDF:
\[
  E_i = n \cdot \Prob(a_i < X \le b_i) = n[F(b_i) - F(a_i)].
\]

\begin{workedexample}[--- Testing Normality]
100 measurements, grouped:
\begin{center}
\begin{tabular}{lcccccc}
\toprule
Interval  & $<45$ & $45$--$50$ & $50$--$55$ & $55$--$60$ & $60$--$65$ & $>65$ \\
Observed  & 5     & 22          & 38          & 24          & 9           & 2     \\
\bottomrule
\end{tabular}
\end{center}
From the data, $\bar{x}=52.3$ and $s=5.1$.  Test $H_0$: data $\sim N(52.3,5.1^2)$.

\begin{solution}
Compute expected proportions using $Z=(X-52.3)/5.1$:
\begin{center}
\begin{tabular}{lccrc}
\toprule
Interval & $z_{\text{lower}}$ & $z_{\text{upper}}$ & $E$ & $(O-E)^2/E$ \\
\midrule
$<45$         & $-\infty$ & $-1.43$ & $100\times0.0764=7.64$ & 0.91 \\
$45$--$50$    & $-1.43$ & $-0.45$ & $100\times0.2490=24.90$ & 0.34 \\
$50$--$55$    & $-0.45$ & $0.53$ & $100\times0.3736=37.36$ & 0.011 \\
$55$--$60$    & $0.53$ & $1.51$ & $100\times0.2404=24.04$ & 0.000 \\
$60$--$65$    & $1.51$ & $2.49$ & $100\times0.0601=6.01$ & 1.49 \\
$>65$         & $2.49$ & $\infty$ & $100\times0.0064=0.64$ & merge \\
\bottomrule
\end{tabular}
\end{center}
Merge last two: $O=11$, $E=6.65$, $(11-6.65)^2/6.65=2.85$.

$X^2=0.91+0.34+0.011+0.000+2.85=4.111$.

$\nu = 5 - 1 - 2 = 2$ (5 cells after merging; 2 parameters $\mu,\sigma$ estimated).

$\chi^2_{0.05,2}=5.991$.  Since $4.111<5.991$, do not reject $H_0$.
\end{solution}
\end{workedexample}

\section{Testing Association: Contingency Tables}
\label{sec:contingency}

\subsection{Setting}

We have two categorical variables $A$ (with $r$ levels) and $B$ (with $c$
levels).  Each of $n$ individuals is classified by both variables, giving an
$r\times c$ \textbf{contingency table} of observed frequencies $O_{ij}$.

\[
  H_0: A \text{ and } B \text{ are independent}
  \quad\text{vs}\quad
  H_1: A \text{ and } B \text{ are associated.}
\]

\subsection{Expected Frequencies Under Independence}

Under $H_0$, $\Prob(\text{row }i, \text{col }j) = \Prob(\text{row }i)\times\Prob(\text{col }j)$.
Estimating marginal probabilities from the data:
\[
  E_{ij} = \frac{R_i \times C_j}{n},
\]
where $R_i$ is the $i$-th row total and $C_j$ is the $j$-th column total.

The test statistic is $X^2 = \sum_{i,j}\frac{(O_{ij}-E_{ij})^2}{E_{ij}}$,
which follows $\chi^2((r-1)(c-1))$ under $H_0$.

\begin{intuition}
We lose one degree of freedom per row margin and one per column margin.
With $r$ row categories and $c$ column categories, we estimate $r+c-1$
parameters from the margins (the $-1$ accounts for the constraint that all
proportions sum to 1), leaving $(r-1)(c-1)$ free cells.
\end{intuition}

\begin{workedexample}[--- Contingency Table Test]
Is blood type associated with hair colour?  Sample of 200 people:
\begin{center}
\begin{tabular}{lcccc|c}
\toprule
           & Blonde & Brown & Black & Red & Total \\
\midrule
Type A     & 22     & 38    & 14    & 6   & 80 \\
Type B     & 12     & 18    & 10    & 5   & 45 \\
Type O     & 28     & 36    & 22    & 8   & 94 \\[2pt]
Type AB    & 4      & 11    & 3     & 3   & 21 \\
\midrule
Total      & 66     & 103   & 49    & 22  & 240 \\
\bottomrule
\end{tabular}
\end{center}
Test at 5\% significance.

\begin{solution}
$H_0$: blood type and hair colour are independent.

Expected frequencies $E_{ij}=R_i C_j/240$:
\begin{center}
\begin{tabular}{lcccc}
\toprule
  & Blonde & Brown & Black & Red \\
\midrule
A  & $80\times66/240=22.0$ & $34.3$ & $16.3$ & $7.3$ \\
B  & $12.4$ & $19.3$ & $9.2$ & $4.1$ \\
O  & $25.9$ & $40.4$ & $19.2$ & $8.6$ \\
AB & $5.8$ & $9.0$ & $4.3$ & $1.9$ \\
\bottomrule
\end{tabular}
\end{center}

Check all $E_{ij}\ge 5$: the cell AB/Red has $E=1.9<5$.  Merge Red into
Black to give a $4\times3$ table and recompute $E$ values.

After merging (Black+Red):
\[
  X^2 \approx 5.12 \quad \text{(detailed calculation omitted)}.
\]
$\nu=(4-1)(3-1)=6$.  $\chi^2_{0.05,6}=12.59$.  Since $5.12<12.59$, we do
not reject $H_0$.  No significant association between blood type and hair colour.
\end{solution}
\end{workedexample}

\begin{realworld}[: Epidemiology]
Contingency table tests are the workhorses of epidemiology.  The landmark
1950 study by Doll and Hill used a $2\times2$ contingency table to establish
the association between smoking and lung cancer---one of the most consequential
statistical analyses in history.
\end{realworld}

\section{Common Mistakes in Chapter 10}

\begin{warningbox}
\begin{enumerate}[label=\textbf{\arabic*.}]
  \item \textbf{Using raw data instead of frequencies.}  The $\chi^2$ test
        uses frequencies (counts), not individual measurements.
  \item \textbf{Not checking $E\ge5$.}  If any $E<5$, merge categories.
  \item \textbf{Wrong degrees of freedom.}  Remember: $\nu=k-1-m$ for
        goodness of fit (where $m$ is the number of estimated parameters);
        $\nu=(r-1)(c-1)$ for contingency tables.
  \item \textbf{Concluding causation from association.}  A significant
        $\chi^2$ test shows association, not causation.
  \item \textbf{Forgetting to estimate parameters from data.}  If you
        fit a Poisson with $\hat{\lambda}=\bar{x}$, subtract 1 from $\nu$.
\end{enumerate}
\end{warningbox}

\section{Summary of Chapter 10}

\begin{summarybox}
\begin{description}[style=nextline]
  \item[Test statistic] $X^2=\sum(O-E)^2/E$; approximately $\chi^2(\nu)$ under $H_0$.
  \item[Goodness of fit] $\nu=k-1-m$; $E_i=np_i$; require $E_i\ge5$.
  \item[Continuous GoF] Group data into intervals; $E_i=n[F(b_i)-F(a_i)]$.
  \item[Contingency tables] $\nu=(r-1)(c-1)$; $E_{ij}=R_iC_j/n$.
  \item[Decision rule] Reject $H_0$ if $X^2>\chi^2_{\alpha,\nu}$ (right-tailed).
  \item[Caution] Association $\ne$ causation; always check expected frequencies.
\end{description}
\end{summarybox}

% ══════════════════════════════════════════════════════════
%  CHAPTER 11: NON-PARAMETRIC TESTS
% ══════════════════════════════════════════════════════════
\chapter{Non-Parametric Tests}
\label{ch:nonpar}

\section{Introduction}

The $t$-tests of Chapter 9 assume that the populations are \emph{normally
distributed}.  What do we do when:
\begin{itemize}
  \item The data are heavily skewed?
  \item The sample is very small and we cannot assess normality?
  \item The measurements are on an \textbf{ordinal scale} (rankings) rather
        than interval/ratio scale?
  \item We have outliers that distort the mean?
\end{itemize}

\textbf{Non-parametric tests} (also called \emph{distribution-free tests})
make fewer assumptions about the underlying distribution.  Instead of means,
they typically use \textbf{ranks} or \textbf{signs}---features that are
robust to outliers and non-normality.

\begin{historicalbox}
The sign test is one of the oldest statistical tests, used informally since
antiquity and formalised in the early 20th century.  Frank Wilcoxon (1892--1965)
proposed the signed-rank and rank-sum tests in 1945, revolutionising
non-parametric statistics.  His methods are now among the most widely used
in medicine and psychology.
\end{historicalbox}

\section{Overview of Non-Parametric Tests}
\label{sec:nonpar-overview}

\begin{center}
\begin{tabular}{lll}
\toprule
\textbf{Test} & \textbf{Setting} & \textbf{Tests} \\
\midrule
Sign test (single-sample)            & One sample            & Median = $m_0$ \\
Wilcoxon signed-rank (single)        & One sample            & Median = $m_0$ (symmetric) \\
Sign test (paired)                   & Paired data           & No difference \\
Wilcoxon matched-pairs signed-rank   & Paired data (symmetric) & No difference \\
Wilcoxon rank-sum                    & Two independent samples & Equal location \\
\bottomrule
\end{tabular}
\end{center}

\section{The Single-Sample Sign Test}
\label{sec:sign-single}

\subsection{Setting}

We wish to test whether the population \textbf{median} equals some
hypothesised value $m_0$.

\[
  H_0: m = m_0 \quad\text{vs}\quad H_1: m \ne m_0 \text{ (or one-sided)}.
\]

\subsection{Procedure}

\begin{enumerate}[label=\textbf{Step \arabic*:}]
  \item For each observation $x_i$, record the \textbf{sign} of $x_i - m_0$:
        a $+$ if $x_i > m_0$, a $-$ if $x_i < m_0$.  Discard ties ($x_i=m_0$).
  \item Let $n^+$ = number of $+$ signs and $n$ = total signs used.
  \item Under $H_0$, each sign is $+$ with probability $1/2$ (since half the
        population lies above the median).  So the number of $+$ signs,
        $S\sim B(n,1/2)$ under $H_0$.
  \item Find the $p$-value using the binomial distribution.
\end{enumerate}

\begin{workedexample}[--- Single-Sample Sign Test]
A nutritionist claims the median daily calorie intake for adults in a region
is 2000 kcal.  A random sample of 12 adults gave the following:
\[
  1850, 2150, 1780, 2300, 1950, 2050, 1700, 2200, 1900, 2400, 2100, 1800.
\]
Test $H_0: m=2000$ vs $H_1: m\ne2000$ at 5\%.

\begin{solution}
Signs of $x_i - 2000$:
$-,+,-,+,-,+,-,+,-,+,+,-$ $\longrightarrow$ 5 positive, 7 negative, $n=12$.

$S \sim B(12, 1/2)$ under $H_0$.  Two-tailed test: find
$\Prob(S\le5)+\Prob(S\ge7)$.

By symmetry of $B(12,1/2)$:
$\Prob(S\le5)=\Prob(S\ge7)$.

$\Prob(S\le5) = \sum_{k=0}^5\binom{12}{k}(1/2)^{12}$
$=\frac{1+12+66+220+495+792}{4096}=\frac{1586}{4096}\approx0.387$.

$p$-value $= 2\times0.387=0.774$.

Since $0.774>0.05$, do not reject $H_0$.  No significant evidence that the
median differs from 2000 kcal.
\end{solution}
\end{workedexample}

\section{The Single-Sample Wilcoxon Signed-Rank Test}
\label{sec:wilcoxon-single}

\subsection{Motivation and Setting}

The sign test only uses the \emph{direction} of deviation from $m_0$,
discarding information about \emph{magnitude}.  The Wilcoxon signed-rank
test is more powerful because it also uses the sizes of deviations.

\textbf{Additional assumption:} The population distribution is
\textbf{symmetric} about its median.  (This is weaker than normality.)

\[
  H_0: m = m_0 \quad\text{vs}\quad H_1: m \ne m_0 \text{ (or one-sided)}.
\]

\subsection{Procedure}

\begin{enumerate}[label=\textbf{Step \arabic*:}]
  \item Compute $d_i = x_i - m_0$.  Discard any $d_i=0$.
  \item Rank the $|d_i|$ from smallest (rank 1) to largest (rank $n$).
        If there are \textbf{ties}, assign the \textbf{average rank}.
  \item Assign each rank a $+$ or $-$ according to the sign of $d_i$.
  \item Compute $T^+$ = sum of positive ranks and $T^-$ = sum of negative ranks.
        Note $T^+ + T^- = n(n+1)/2$.
  \item The test statistic is $T = \min(T^+, T^-)$ (for two-tailed test).
  \item Reject $H_0$ if $T \le T_{\text{crit}}$ (from tables; small $T$
        gives evidence against $H_0$).
\end{enumerate}

\begin{intuition}
Under $H_0$, positive and negative deviations should be roughly balanced in
both number and size.  $T^+$ and $T^-$ should be approximately equal.  If
$T^+$ is much larger than $T^-$ (or vice versa), the distribution appears to
be shifted away from $m_0$.
\end{intuition}

\begin{workedexample}[--- Single-Sample Wilcoxon Test]
Using the data from the previous example
($1850,2150,1780,2300,1950,2050,1700,2200,1900,2400,2100,1800$),
test $H_0:m=2000$ vs $H_1:m\ne2000$ using the Wilcoxon signed-rank test
at 5\%.

\begin{solution}
Compute $d_i = x_i-2000$ and rank $|d_i|$:

\begin{center}
\begin{tabular}{rrrrl}
\toprule
$x_i$ & $d_i$ & $|d_i|$ & Rank & Signed rank \\
\midrule
1700 & $-300$ & 300 & 12 & $-12$ \\
1780 & $-220$ & 220 & 10 & $-10$ \\
1800 & $-200$ & 200 & 9  & $-9$  \\
1850 & $-150$ & 150 & 7  & $-7$  \\
1900 & $-100$ & 100 & 4  & $-4$  \\
1950 & $-50$  & 50  & 2  & $-2$  \\
2050 & $+50$  & 50  & 2  & $+2$  \\
2100 & $+100$ & 100 & 4  & $+4$  \\
2150 & $+150$ & 150 & 7  & $+7$  \\
2200 & $+200$ & 200 & 9  & $+9$  \\
2300 & $+300$ & 300 & 12 & $+12$ \\
2400 & $+400$ & 400 & 13 & $+13$ \\
\bottomrule
\end{tabular}
\end{center}

Note: $|d|=50$ appears twice (tied), so both get rank $(1+2)/2=1.5$;
similarly $|d|=100$ gives rank $(3+4)/2=3.5$; $|d|=150$ gives rank $5.5$;
$|d|=200$ gives rank $7.5$; $|d|=300$ gives rank $10.5$.

Recomputed:
$T^- = 12+10.5+7.5+5.5+3.5+1.5 = 40.5$.
$T^+ = 1.5+3.5+5.5+7.5+10.5+13 = 41.5$.
Check: $40.5+41.5=82=12\times13/2$.\ \checkmark

$T=\min(40.5,41.5)=40.5$.

From Wilcoxon tables with $n=12$, $\alpha=0.05$ two-tailed: $T_{\text{crit}}=13$.

Since $40.5 > 13$, do not reject $H_0$.  Consistent with the sign test result.
\end{solution}
\end{workedexample}

\section{The Paired-Sample Sign Test}
\label{sec:sign-paired}

For paired data $(x_{1i},x_{2i})$, compute $d_i=x_{1i}-x_{2i}$ and apply
the single-sample sign test to the $d_i$ with $m_0=0$.

\[
  H_0: m_D = 0 \quad\text{vs}\quad H_1: m_D\ne0.
\]

This tests whether the median difference is zero---i.e., whether the two
treatments tend to give the same result.

\section{The Wilcoxon Matched-Pairs Signed-Rank Test}
\label{sec:wilcoxon-paired}

\subsection{Setting}

Paired data $(x_{1i},x_{2i})$, $i=1,\ldots,n$.  Population symmetric.
$H_0$: no difference in location; $H_1$: there is a difference.

Procedure: compute $d_i=x_{1i}-x_{2i}$, then apply the single-sample
Wilcoxon signed-rank test to the $d_i$ with $m_0=0$.

\begin{workedexample}[--- Wilcoxon Matched-Pairs Test]
Eight patients were given two analgesics (A and B) on different occasions.
Pain scores (out of 10, lower is better):
\begin{center}
\begin{tabular}{lrrrrrrrr}
\toprule
Patient & 1 & 2 & 3 & 4 & 5 & 6 & 7 & 8 \\
\midrule
Drug A & 6.2 & 4.8 & 7.1 & 5.5 & 6.8 & 4.2 & 7.5 & 5.0 \\
Drug B & 5.5 & 4.5 & 6.8 & 5.0 & 7.2 & 3.8 & 6.9 & 4.7 \\
$d=A-B$ & 0.7 & 0.3 & 0.3 & 0.5 & $-0.4$ & 0.4 & 0.6 & 0.3 \\
\bottomrule
\end{tabular}
\end{center}
Test at 5\% (two-tailed) whether there is a difference between drugs.

\begin{solution}
Ranks of $|d_i|$:
\begin{center}
\begin{tabular}{cccl}
\toprule
$|d_i|$ & Count & Rank & Signed rank \\
\midrule
0.3 & 3 & $(1+2+3)/3=2$ & $+2,+2,+2$ \\
0.4 & 2 & $(4+5)/2=4.5$ & $-4.5,+4.5$ \\
0.5 & 1 & 6 & $+6$ \\
0.6 & 1 & 7 & $+7$ \\
0.7 & 1 & 8 & $+8$ \\
\bottomrule
\end{tabular}
\end{center}

$T^+ = 2+2+2+4.5+6+7+8 = 31.5$.
$T^- = 4.5$.
$T = \min(31.5, 4.5) = 4.5$.

From tables, $n=8$, $\alpha=0.05$ two-tailed: $T_{\text{crit}}=4$.

Since $4.5 > 4$, we do not reject $H_0$ at 5\%.
(The result is borderline; at 10\% the critical value is 6 and we would
reject.)
\end{solution}
\end{workedexample}

\section{The Wilcoxon Rank-Sum Test}
\label{sec:rank-sum}

\subsection{Setting}

Two \textbf{independent} samples, sizes $n_1\le n_2$.  No normality assumed.
$H_0$: populations have the same distribution (or same location).

\subsection{Procedure}

\begin{enumerate}[label=\textbf{Step \arabic*:}]
  \item Combine both samples into one list of $n_1+n_2$ observations and
        rank all values from 1 to $n_1+n_2$ (handling ties with average ranks).
  \item Compute $W$ = sum of ranks from \textbf{sample 1} (the smaller group).
  \item Compare $W$ with critical values from tables.
        Reject $H_0$ if $W\le W_{\text{lower}}$ or $W\ge W_{\text{upper}}$
        (two-tailed).
\end{enumerate}

\begin{intuition}
If $H_0$ is true and the populations have the same location, the ranks should
be roughly evenly distributed between the two samples.  A very small or very
large rank sum for sample 1 suggests that sample 1 tends to come from a
lower or higher distribution.
\end{intuition}

\begin{workedexample}[--- Wilcoxon Rank-Sum Test]
Compare the daily earnings (£) of two groups of freelancers:
\[
  \text{Group A (}n_1=6\text{):}\quad 85, 112, 97, 134, 76, 108.
  \quad
  \text{Group B (}n_2=7\text{):}\quad 95, 128, 143, 89, 117, 104, 136.
\]
Test at 5\% (two-tailed) whether the populations have equal location.

\begin{solution}
Combined and ranked:
\begin{center}
\begin{tabular}{crcc}
\toprule
Value & Rank & Group & \\
\midrule
76  & 1  & A & \\
85  & 2  & A & \\
89  & 3  & B & \\
95  & 4  & B & \\
97  & 5  & A & \\
104 & 6  & B & \\
108 & 7  & A & \\
112 & 8  & A & \\
117 & 9  & B & \\
128 & 10 & B & \\
134 & 11 & A & \\
136 & 12 & B & \\
143 & 13 & B & \\
\bottomrule
\end{tabular}
\end{center}

$W$ = sum of Group A ranks $= 1+2+5+7+8+11 = 34$.

From tables with $n_1=6$, $n_2=7$ at 5\% two-tailed:
$W_{\text{lower}}=28$, $W_{\text{upper}}=56$.

Since $28 \le 34 \le 56$, do not reject $H_0$.  No significant evidence of
a difference in earnings between the two groups.
\end{solution}
\end{workedexample}

\subsection{Large Sample Approximation}

For large $n_1,n_2$, under $H_0$:
\[
  \E(W) = \frac{n_1(n_1+n_2+1)}{2}, \qquad
  \Var(W) = \frac{n_1 n_2(n_1+n_2+1)}{12}.
\]
Use $Z = \frac{W - \E(W)}{\sqrt{\Var(W)}} \approx N(0,1)$.

\section{Comparing Parametric and Non-Parametric Tests}

\begin{center}
\begin{tabular}{lll}
\toprule
\textbf{Aspect} & \textbf{Parametric} & \textbf{Non-parametric} \\
\midrule
Assumptions     & Normality (approximately) & Minimal \\
Measure tested  & Mean                & Median / Rank \\
Power           & Higher (if assumptions met) & Lower \\
Robustness      & Sensitive to outliers & Robust \\
Data type       & Interval/ratio      & Ordinal or above \\
Sample size     & More important      & Usable with small $n$ \\
\bottomrule
\end{tabular}
\end{center}

\begin{realworld}[: Psychology and Medicine]
Non-parametric tests are the standard in studies with small samples, Likert
scale data (e.g., ``rate your pain from 1 to 10''), or non-normally
distributed outcomes.  Clinical trials with fewer than 30 participants often
use the Wilcoxon tests.  The rank-sum test also underlies the Mann--Whitney U
test, widely used in social science research.
\end{realworld}

\section{Common Mistakes in Chapter 11}

\begin{warningbox}
\begin{enumerate}[label=\textbf{\arabic*.}]
  \item \textbf{Using the wrong test for the data structure.}  Matched-pairs
        for paired data; rank-sum for independent samples.
  \item \textbf{Mishandling ties.}  Always use average ranks for tied values.
  \item \textbf{Incorrect critical value direction.}  For signed-rank and
        rank-sum tests, we reject when $T$ (or $W$) is \emph{small} (two-tailed:
        outside the critical region defined by $[W_L, W_U]$).
  \item \textbf{Discarding ties incorrectly.}  For sign tests, discard
        $d_i=0$.  For Wilcoxon signed-rank, discard $d_i=0$.
  \item \textbf{Forgetting that non-parametric tests have less power.}
        If normality is plausible, the $t$-test is preferable.
\end{enumerate}
\end{warningbox}

\section{Summary of Chapter 11}

\begin{summarybox}
\begin{description}[style=nextline]
  \item[Sign test] Uses $+/-$ signs only; $S\sim B(n,1/2)$ under $H_0$;
    tests the median.
  \item[Wilcoxon signed-rank] Uses signed ranks of $|d_i|$; $T=\min(T^+,T^-)$;
    requires symmetry; more powerful than sign test.
  \item[Matched-pairs] Apply signed-rank test to differences $d_i=x_{1i}-x_{2i}$.
  \item[Rank-sum] Combine and rank both samples; use sum of ranks from
    smaller sample $W$; tests for equal location.
  \item[All non-parametric tests] Robust to non-normality; use when normality
    cannot be assumed; lose some power compared to $t$-tests.
\end{description}
\end{summarybox}

% ══════════════════════════════════════════════════════════
%  CHAPTER 12: PROBABILITY GENERATING FUNCTIONS
% ══════════════════════════════════════════════════════════
\chapter{Probability Generating Functions}
\label{ch:pgf}

\section{Introduction}

Imagine you could encode the entire probability distribution of a discrete
random variable into a single function.  From that function, you could read
off probabilities, compute the mean, compute the variance, and determine the
distribution of the sum of independent copies---all by algebraic
manipulation.

This is exactly what the \textbf{probability generating function} (PGF) does.
It is one of the most elegant tools in probability theory, connecting
combinatorics, generating functions, and probability in a profound way.

PGFs are used in:
\begin{itemize}
  \item \textbf{Queueing theory:} analysing waiting times in call centres,
        computer networks, and hospitals.
  \item \textbf{Branching processes:} modelling population growth,
        nuclear chain reactions, and the spread of disease.
  \item \textbf{Actuarial science:} computing aggregate insurance claims.
  \item \textbf{Combinatorics:} counting lattice paths, partitions, and
        tree structures.
  \item \textbf{Algorithm analysis:} average-case complexity of algorithms.
\end{itemize}

\begin{historicalbox}
Generating functions were introduced by Abraham de Moivre (1667--1754) and
developed by Euler, Laplace, and Gauss.  The probability generating function
as a tool for random variables was formalised in the 20th century.  The
moment generating function (a close cousin, defined for both discrete and
continuous RVs) was central to the work of Laplace and later to the
mathematical foundations of statistics.
\end{historicalbox}

\section{The Probability Generating Function}
\label{sec:pgf-def}

\subsection{Definition}

\begin{definition}[Probability Generating Function]
Let $X$ be a \textbf{non-negative integer-valued} random variable with
$\Prob(X=k) = p_k$ for $k=0,1,2,\ldots$.  The \textbf{probability generating
function} (PGF) of $X$ is
\[
  G_X(t) = \E(t^X) = \sum_{k=0}^{\infty} p_k\, t^k
  = p_0 + p_1 t + p_2 t^2 + p_3 t^3 + \cdots
\]
for all $t$ in the radius of convergence.
\end{definition}

\begin{intuition}
The PGF is a \emph{power series whose coefficients are the probabilities}.
The probability of $X=k$ is the coefficient of $t^k$ in $G_X(t)$.  This is
why it is called a \emph{generating} function: it generates the probabilities.

The key insight is that the argument $t$ is not a probability or a
value of $X$---it is a formal variable (though we often evaluate at specific
values of $t$).  The PGF always converges at $t=1$ (the sum of all
probabilities equals 1), and for many distributions it converges for $|t|\le1$.
\end{intuition}

\subsection{Properties at Key Values}

\begin{proposition}
For any PGF $G_X(t)$:
\begin{enumerate}[label=\textbf{(\roman*)}]
  \item $G_X(0) = p_0 = \Prob(X=0)$.
  \item $G_X(1) = \sum_{k=0}^\infty p_k = 1$.
  \item The probability of any particular value: $p_k = \dfrac{G_X^{(k)}(0)}{k!}$,
        where $G_X^{(k)}$ denotes the $k$-th derivative.
\end{enumerate}
\end{proposition}

Property (iii) is exactly the Taylor series statement: the coefficient of $t^k$
is $G_X^{(k)}(0)/k!$.

\subsection{PGFs of Common Distributions}

\begin{workedexample}[--- PGF of the Bernoulli Distribution]
$X\sim\operatorname{Bernoulli}(p)$: $P(X=0)=q=1-p$, $P(X=1)=p$.

\begin{solution}
$G_X(t) = p_0\cdot t^0 + p_1\cdot t^1 = q + pt$.
\end{solution}
\end{workedexample}

\begin{workedexample}[--- PGF of the Binomial Distribution]
$X\sim B(n,p)$: $P(X=k)=\binom{n}{k}p^k q^{n-k}$.

\begin{solution}
\[
  G_X(t) = \sum_{k=0}^n\binom{n}{k}p^k q^{n-k}t^k
  = \sum_{k=0}^n\binom{n}{k}(pt)^k q^{n-k}
  = (q+pt)^n.
\]
by the Binomial Theorem.  This is a polynomial of degree $n$---finite!
\end{solution}
\end{workedexample}

\begin{workedexample}[--- PGF of the Geometric Distribution]
$X\sim\operatorname{Geom}(p)$: $P(X=k)=pq^{k-1}$, $k=1,2,\ldots$ (number
of trials until first success).

\begin{solution}
\[
  G_X(t) = \sum_{k=1}^\infty pq^{k-1}t^k = pt\sum_{k=1}^\infty(qt)^{k-1}
  = pt\cdot\frac{1}{1-qt} = \frac{pt}{1-qt},\quad |qt|<1.
\]
\end{solution}
\end{workedexample}

\begin{workedexample}[--- PGF of the Poisson Distribution]
$X\sim\operatorname{Poisson}(\lambda)$: $P(X=k)=e^{-\lambda}\lambda^k/k!$.

\begin{solution}
\[
  G_X(t) = \sum_{k=0}^\infty\frac{e^{-\lambda}\lambda^k}{k!}t^k
  = e^{-\lambda}\sum_{k=0}^\infty\frac{(\lambda t)^k}{k!}
  = e^{-\lambda}\cdot e^{\lambda t}
  = e^{\lambda(t-1)}.
\]
\end{solution}
\end{workedexample}

\section{Mean and Variance from the PGF}
\label{sec:pgf-moments}

\subsection{Derivation}

\begin{theorem}[Moments from the PGF]
\label{thm:pgf-moments}
If $G_X(t) = \E(t^X)$, then:
\begin{enumerate}[label=\textbf{(\roman*)}]
  \item $\E(X) = G_X'(1)$.
  \item $\Var(X) = G_X''(1) + G_X'(1) - [G_X'(1)]^2$.
\end{enumerate}
\end{theorem}

\begin{proof}
Differentiate $G_X(t) = \sum_{k=0}^\infty p_k t^k$ term by term
(valid for $|t|<1$ and by continuity at $t=1$ for distributions with
finite mean):
\[
  G_X'(t) = \sum_{k=1}^\infty k p_k t^{k-1}.
\]
Setting $t=1$:
\[
  G_X'(1) = \sum_{k=1}^\infty k p_k = \E(X). \quad\checkmark
\]

Differentiate again:
\[
  G_X''(t) = \sum_{k=2}^\infty k(k-1)p_k t^{k-2}.
\]
Setting $t=1$:
\[
  G_X''(1) = \sum_{k=2}^\infty k(k-1)p_k = \E[X(X-1)] = \E(X^2) - \E(X).
\]
Therefore $\E(X^2) = G_X''(1) + G_X'(1)$, and
\[
  \Var(X) = \E(X^2) - [\E(X)]^2 = G_X''(1) + G_X'(1) - [G_X'(1)]^2. \quad\checkmark
\]
\end{proof}

\begin{workedexample}[--- Recovering Binomial Moments]
For $X\sim B(n,p)$, $G_X(t)=(q+pt)^n$.  Verify $\E(X)=np$ and
$\Var(X)=npq$.

\begin{solution}
$G_X'(t) = np(q+pt)^{n-1}$, so $G_X'(1) = np(q+p)^{n-1} = np$.\ \checkmark

$G_X''(t) = n(n-1)p^2(q+pt)^{n-2}$, so $G_X''(1) = n(n-1)p^2$.

$\Var(X) = n(n-1)p^2 + np - (np)^2 = n^2p^2 - np^2 + np - n^2p^2
= np(1-p) = npq$.\ \checkmark
\end{solution}
\end{workedexample}

\begin{workedexample}[--- Poisson Moments via PGF]
For $X\sim\operatorname{Poisson}(\lambda)$, $G_X(t)=e^{\lambda(t-1)}$.
Find $\E(X)$ and $\Var(X)$.

\begin{solution}
$G_X'(t) = \lambda e^{\lambda(t-1)}$, so $G_X'(1) = \lambda = \E(X)$.\ \checkmark

$G_X''(t) = \lambda^2 e^{\lambda(t-1)}$, so $G_X''(1) = \lambda^2$.

$\Var(X) = \lambda^2 + \lambda - \lambda^2 = \lambda$.\ \checkmark

Both the mean and variance of the Poisson distribution equal $\lambda$---a
defining property.
\end{solution}
\end{workedexample}

\section{The Sum of Independent Random Variables}
\label{sec:pgf-sum}

\subsection{The Key Theorem}

The most powerful property of the PGF is how it handles sums.

\begin{theorem}[PGF of a Sum]
\label{thm:pgf-sum}
Let $X$ and $Y$ be \textbf{independent} non-negative integer-valued random
variables.  Then
\[
  G_{X+Y}(t) = G_X(t)\cdot G_Y(t).
\]
\end{theorem}

\begin{proof}
By independence of $X$ and $Y$:
\[
  G_{X+Y}(t) = \E(t^{X+Y}) = \E(t^X\cdot t^Y) = \E(t^X)\cdot\E(t^Y)
  = G_X(t)\cdot G_Y(t).
\]
The factorisation uses independence: $\E(UV)=\E(U)\E(V)$ when $U$ and $V$
are independent.
\end{proof}

\begin{intuition}
Multiplying power series corresponds to \textbf{convolving} the coefficient
sequences.  The coefficient of $t^n$ in $G_X(t)\cdot G_Y(t)$ is
$\sum_{k=0}^n p_k q_{n-k}$, where $p_k=P(X=k)$ and $q_j=P(Y=j)$.  But this
is exactly $P(X+Y=n)=\sum_{k=0}^n P(X=k)P(Y=n-k)$ (the convolution formula
for sums of independent RVs).  PGFs \emph{turn convolution into
multiplication}---a massive simplification.
\end{intuition}

\begin{corollary}[Sum of $n$ IID variables]
If $X_1,X_2,\ldots,X_n$ are independent, each with PGF $G(t)$, then
\[
  G_{X_1+X_2+\cdots+X_n}(t) = [G(t)]^n.
\]
\end{corollary}

\subsection{Applications}

\begin{workedexample}[--- Sum of Poisson Variables]
Let $X\sim\operatorname{Poisson}(\lambda)$ and $Y\sim\operatorname{Poisson}(\mu)$
be independent.  Show that $X+Y\sim\operatorname{Poisson}(\lambda+\mu)$.

\begin{solution}
\[
  G_{X+Y}(t) = G_X(t)\cdot G_Y(t) = e^{\lambda(t-1)}\cdot e^{\mu(t-1)}
  = e^{(\lambda+\mu)(t-1)}.
\]
This is the PGF of a $\operatorname{Poisson}(\lambda+\mu)$ distribution.
Since the PGF uniquely determines the distribution, $X+Y\sim\operatorname{Poisson}(\lambda+\mu)$.
\end{solution}
\end{workedexample}

\begin{workedexample}[--- Sum of Binomial Variables]
If $X\sim B(m,p)$ and $Y\sim B(n,p)$ independently, find the distribution of $X+Y$.

\begin{solution}
\[
  G_{X+Y}(t)=(q+pt)^m\cdot(q+pt)^n=(q+pt)^{m+n}.
\]
This is the PGF of $B(m+n,p)$.  Hence $X+Y\sim B(m+n,p)$.
\end{solution}
\end{workedexample}

\begin{workedexample}[--- Non-standard Application]
A bus company operates two independent routes.  On Route 1, the number of
passengers per journey follows $\operatorname{Poisson}(8)$.  On Route 2, it
follows $\operatorname{Poisson}(5)$.  What is the probability that the total
number of passengers on both routes in one journey is exactly 10?

\begin{solution}
Total $T = X+Y\sim\operatorname{Poisson}(13)$.
\[
  P(T=10) = e^{-13}\cdot\frac{13^{10}}{10!}
  = e^{-13}\cdot\frac{137858491849}{3628800}
  \approx 2.203\times10^{-6}\times 37974.6
  \approx 0.0853.
\]
\end{solution}
\end{workedexample}

\section{Three or More Random Variables}
\label{sec:pgf-three}

The result extends immediately to three or more independent variables:
\[
  G_{X_1+\cdots+X_n}(t) = G_{X_1}(t)\cdot G_{X_2}(t)\cdots G_{X_n}(t).
\]

When the variables are not identically distributed, we simply multiply
\emph{different} PGFs.

\begin{workedexample}[--- Mixed Sum]
$X\sim\operatorname{Poisson}(3)$, $Y\sim B(4,0.5)$, and $Z\sim\operatorname{Geom}(0.4)$
are independent.  Find $\E(X+Y+Z)$ and $\Var(X+Y+Z)$.

\begin{solution}
By linearity:
\[
  \E(X+Y+Z) = \E(X)+\E(Y)+\E(Z) = 3 + 4(0.5) + \frac{1}{0.4} = 3+2+2.5 = 7.5.
\]

By independence:
\[
  \Var(X+Y+Z) = \Var(X)+\Var(Y)+\Var(Z) = 3 + 4(0.5)(0.5) + \frac{0.6}{0.16}
  = 3 + 1 + 3.75 = 7.75.
\]

Alternatively, using PGFs:
$G_{X+Y+Z}(t) = e^{3(t-1)}\cdot(0.5+0.5t)^4\cdot\frac{0.4t}{1-0.6t}$.
Differentiating at $t=1$ gives the same results.
\end{solution}
\end{workedexample}

\subsection{The Random Sum}

A particularly beautiful application is the \textbf{random sum}.  Suppose
$N$ is a random variable and, given $N=n$, we sum $n$ independent copies
of $X$.  The total $S = X_1+X_2+\cdots+X_N$ is a random sum.

\begin{theorem}[PGF of a Random Sum]
If $N$ and $X_1,X_2,\ldots$ are independent and $X_i$ are iid with PGF
$G_X(t)$, and $N$ has PGF $G_N(t)$, then
\[
  G_S(t) = G_N(G_X(t)).
\]
\end{theorem}

\begin{proof}
Condition on $N$:
\[
  G_S(t) = \E(t^S) = \sum_{n=0}^\infty\E(t^S\mid N=n)\Prob(N=n)
  = \sum_{n=0}^\infty[G_X(t)]^n\Prob(N=n)
  = G_N(G_X(t)).
\]
\end{proof}

\begin{realworld}[: Actuarial Science]
The random sum models \emph{aggregate claims}: $N$ is the (random) number of
claims in a period, and $X_i$ is the (random) amount of the $i$-th claim.
The total loss $S=\sum_{i=1}^N X_i$ is exactly the random sum.  PGFs (and
their continuous analogues, moment generating functions and Laplace
transforms) are used to compute the distribution of $S$ and hence to
price insurance products.
\end{realworld}

\begin{workedexample}[--- Random Sum]
Customers arrive at a shop according to a $\operatorname{Poisson}(5)$ process
(so the number of customers $N$ per hour has $\operatorname{Poisson}(5)$
distribution).  Each customer independently spends a number of items $X$
with $P(X=1)=0.6$ and $P(X=2)=0.4$ (so $X$ takes values 1 or 2).  Find
$\E(S)$ and $\Var(S)$ where $S$ is the total number of items sold per hour.

\begin{solution}
$G_N(t)=e^{5(t-1)}$ and $G_X(t)=0.6t+0.4t^2$.

$\E(N)=5$, $\Var(N)=5$.
$\E(X) = 0.6+0.8=1.4$, $\E(X^2)=0.6+1.6=2.2$, $\Var(X)=2.2-1.96=0.24$.

By the \textbf{Wald's identity}:
\[
  \E(S) = \E(N)\cdot\E(X) = 5\times1.4 = 7.
\]
\[
  \Var(S) = \E(N)\cdot\Var(X) + \Var(N)\cdot[\E(X)]^2
  = 5(0.24) + 5(1.96) = 1.2 + 9.8 = 11.
\]
\end{solution}
\end{workedexample}

\section{Uniqueness of the PGF}

\begin{theorem}[Uniqueness]
Two discrete non-negative integer-valued random variables have the same
distribution if and only if they have the same PGF.
\end{theorem}

This is crucial: it means that if you recognise the PGF of $X+Y$ as the
PGF of a known distribution, you can immediately conclude $X+Y$ has that
distribution---as we did in the Poisson and binomial examples above.

\section{Connection to Moment Generating Functions}

\begin{remark}
The \textbf{moment generating function} (MGF) of $X$ is defined as
$M_X(s) = \E(e^{sX})$.  For a non-negative integer-valued $X$:
\[
  G_X(t) = M_X(\ln t), \qquad M_X(s) = G_X(e^s).
\]
The MGF is defined for all random variables (not just non-negative
integer-valued ones) and is the tool used at university level for proving
the Central Limit Theorem and working with continuous distributions.
\end{remark}

\section{Advanced Exploration: Branching Processes}

\begin{exploration}
\textbf{Extinction of Surnames and Species}

A \textbf{Galton--Watson branching process} models a population where each
individual in generation $n$ independently produces a random number of
offspring $X$ with PGF $G(t)$.  The population size in generation $n+1$ is
the sum of all offspring.

\begin{enumerate}
  \item Show that the PGF of the generation-$n$ population size is the
        $n$-fold composition: $G_n(t) = G(G(\cdots G(t)\cdots))$.
  \item Let $q$ be the probability of eventual extinction (the population
        size reaches 0).  Show that $q$ satisfies $G(q)=q$.
  \item Show that if $\E(X)=G'(1)\le 1$, then $q=1$ (extinction is
        certain).  If $\E(X)>1$, then $q<1$.
  \item This was used by Francis Galton (1822--1911) to study the extinction
        of surnames in Victorian England.  Research the fascinating history
        of this problem.
\end{enumerate}
\end{exploration}

\section{Common Mistakes in Chapter 12}

\begin{warningbox}
\begin{enumerate}[label=\textbf{\arabic*.}]
  \item \textbf{Evaluating $G_X(1)$ and expecting a useful probability.}
        $G_X(1)=1$ always; this just confirms normalisation.
  \item \textbf{Computing $\E(X) = G_X'(t)$ without setting $t=1$.}
        Remember: $\E(X) = G_X'(1)$, not $G_X'(t)$.
  \item \textbf{Forgetting independence} when multiplying PGFs.
        $G_{X+Y}(t)=G_X(t)G_Y(t)$ \emph{only if} $X$ and $Y$ are independent.
  \item \textbf{Incorrect formula for variance.}  Use
        $\Var(X)=G''(1)+G'(1)-[G'(1)]^2$, not just $G''(1)$.
  \item \textbf{Confusing PGF with MGF.}  The PGF is defined only for
        non-negative integer-valued RVs; the MGF is more general.
\end{enumerate}
\end{warningbox}

\section{Summary of Chapter 12}

\begin{summarybox}
\begin{description}[style=nextline]
  \item[Definition] $G_X(t)=\E(t^X)=\sum_{k=0}^\infty p_k t^k$.
    Coefficient of $t^k$ gives $P(X=k)$.
  \item[Key values] $G(0)=P(X=0)$; $G(1)=1$.
  \item[Mean] $\E(X)=G'(1)$.
  \item[Variance] $\Var(X)=G''(1)+G'(1)-[G'(1)]^2$.
  \item[Sum of independent RVs] $G_{X+Y}(t)=G_X(t)\cdot G_Y(t)$.
  \item[IID sum] $G_{S_n}(t)=[G_X(t)]^n$.
  \item[Random sum] $G_S(t)=G_N(G_X(t))$.
  \item[Uniqueness] PGF uniquely determines the distribution.
\end{description}

\medskip
\textbf{PGFs of standard distributions:}
\begin{center}
\begin{tabular}{lll}
\toprule
Distribution & Parameters & $G_X(t)$ \\
\midrule
Bernoulli   & $p$       & $q+pt$ \\
Binomial    & $n,p$     & $(q+pt)^n$ \\
Geometric   & $p$       & $pt/(1-qt)$ \\
Poisson     & $\lambda$ & $e^{\lambda(t-1)}$ \\
\bottomrule
\end{tabular}
\end{center}
\end{summarybox}

% ══════════════════════════════════════════════════════════
%  FINAL REVIEW: EXTENSION PROBLEMS
% ══════════════════════════════════════════════════════════
\chapter*{Extension and Review Problems}
\addcontentsline{toc}{chapter}{Extension and Review Problems}

The following problems range from challenging A-Level questions to
university-entrance level.  They are designed to test deep understanding
and synthesis across chapters.

\section*{Chapter 8: Continuous Random Variables}

\begin{exercise}
Let $X$ have PDF $f(x)=cx^2(1-x)^3$ for $0\le x\le1$.
\begin{enumerate}[label=\textbf{(\alph*)}]
  \item Find $c$ (hint: this is a Beta distribution).
  \item Find $\E(X)$ and $\Var(X)$.
  \item Find the CDF $F(x)$.
  \item Let $Y=-2\ln X$.  Find the PDF of $Y$ and identify its distribution.
\end{enumerate}
\end{exercise}

\begin{exercise}[University level]
Prove that if $X$ has a continuous CDF $F$, then $U=F(X)\sim U(0,1)$.
(This is the \textbf{probability integral transform}.)  How does this
relate to simulation?
\end{exercise}

\section*{Chapter 9: Inferential Statistics}

\begin{exercise}
Two independent samples:
\begin{itemize}
  \item Sample 1: $n_1=8$, $\bar{x}_1=23.4$, $s_1^2=9.6$.
  \item Sample 2: $n_2=10$, $\bar{x}_2=20.1$, $s_2^2=11.2$.
\end{itemize}
\begin{enumerate}[label=\textbf{(\alph*)}]
  \item Test $H_0:\mu_1=\mu_2$ vs $H_1:\mu_1>\mu_2$ at 5\%.
  \item Construct a 95\% CI for $\mu_1-\mu_2$.
  \item If the variances might not be equal, describe what modification
        you would make to the test.
\end{enumerate}
\end{exercise}

\begin{exercise}[Exploration]
Explain intuitively why a paired $t$-test is more powerful than a two-sample
$t$-test when the data are paired.  Use the formula for variance to make
this precise: show that $\Var(D_i)=\Var(X_{1i})+\Var(X_{2i})-2\Cov(X_{1i},X_{2i})$
and interpret what happens when the correlation is positive.
\end{exercise}

\section*{Chapter 10: Chi-Squared Tests}

\begin{exercise}
In a $2\times2$ contingency table with entries $a,b,c,d$ (row totals $a+b$,
$c+d$; column totals $a+c$, $b+d$; grand total $n$), show that
\[
  X^2 = \frac{n(ad-bc)^2}{(a+b)(c+d)(a+c)(b+d)}.
\]
This elegant formula shows that the test statistic depends only on
$|ad-bc|$, the \textbf{cross-product difference}.
\end{exercise}

\section*{Chapter 11: Non-Parametric Tests}

\begin{exercise}
Explain carefully the connection between the Wilcoxon rank-sum test and
the \textbf{Mann--Whitney U statistic} $U$, defined as the number of
pairs $(X_i,Y_j)$ where $X_i>Y_j$.  Show that $W = U + n_1(n_1+1)/2$.
\end{exercise}

\begin{exercise}[University level]
Show that for large $n$, the signed-rank statistic $T^+$ is approximately
normally distributed with mean $n(n+1)/4$ and variance $n(n+1)(2n+1)/24$.
Use this to derive a large-sample $z$-test for the median.
\end{exercise}

\section*{Chapter 12: Probability Generating Functions}

\begin{exercise}
The \textbf{negative binomial} distribution $\operatorname{NB}(r,p)$ gives
the number of trials until the $r$-th success:
$P(X=k) = \binom{k-1}{r-1}p^r q^{k-r}$ for $k=r,r+1,\ldots$.
\begin{enumerate}[label=\textbf{(\alph*)}]
  \item Show that the PGF is $G(t) = \left(\frac{pt}{1-qt}\right)^r$.
  \item Derive $\E(X)=r/p$ and $\Var(X)=rq/p^2$.
  \item Show that the sum of $r$ independent $\operatorname{Geom}(p)$
        variables has the $\operatorname{NB}(r,p)$ distribution.
\end{enumerate}
\end{exercise}

\begin{exercise}[Branching process]
In a Galton--Watson branching process, each individual has 0, 1, or 2
offspring with probabilities $1/3, 1/3, 1/3$ respectively.
\begin{enumerate}[label=\textbf{(\alph*)}]
  \item Find the PGF $G(t)$.
  \item Compute $\E(X)$ and explain what this implies for long-term
        population behaviour.
  \item Find the extinction probability $q$ by solving $G(q)=q$.
\end{enumerate}
\end{exercise}

% ── Appendix ───────────────────────────────────────────────
\appendix
\chapter{Statistical Tables Reference}

\section{Student's $t$-Distribution Critical Values}

\begin{center}
\begin{tabular}{cccccc}
\toprule
$\nu$ & $t_{0.10}$ & $t_{0.05}$ & $t_{0.025}$ & $t_{0.01}$ & $t_{0.005}$ \\
\midrule
1  & 3.078 & 6.314 & 12.706 & 31.821 & 63.657 \\
2  & 1.886 & 2.920 & 4.303  & 6.965  & 9.925  \\
3  & 1.638 & 2.353 & 3.182  & 4.541  & 5.841  \\
4  & 1.533 & 2.132 & 2.776  & 3.747  & 4.604  \\
5  & 1.476 & 2.015 & 2.571  & 3.365  & 4.032  \\
6  & 1.440 & 1.943 & 2.447  & 3.143  & 3.707  \\
7  & 1.415 & 1.895 & 2.365  & 2.998  & 3.499  \\
8  & 1.397 & 1.860 & 2.306  & 2.896  & 3.355  \\
9  & 1.383 & 1.833 & 2.262  & 2.821  & 3.250  \\
10 & 1.372 & 1.812 & 2.228  & 2.764  & 3.169  \\
12 & 1.356 & 1.782 & 2.179  & 2.681  & 3.055  \\
15 & 1.341 & 1.753 & 2.131  & 2.602  & 2.947  \\
20 & 1.325 & 1.725 & 2.086  & 2.528  & 2.845  \\
25 & 1.316 & 1.708 & 2.060  & 2.485  & 2.787  \\
30 & 1.310 & 1.697 & 2.042  & 2.457  & 2.750  \\
40 & 1.303 & 1.684 & 2.021  & 2.423  & 2.704  \\
60 & 1.296 & 1.671 & 2.000  & 2.390  & 2.660  \\
$\infty$ & 1.282 & 1.645 & 1.960 & 2.326 & 2.576 \\
\bottomrule
\end{tabular}
\end{center}

\section{Chi-Squared Critical Values}

\begin{center}
\begin{tabular}{cccccc}
\toprule
$\nu$ & $\chi^2_{0.10}$ & $\chi^2_{0.05}$ & $\chi^2_{0.025}$ & $\chi^2_{0.01}$ & $\chi^2_{0.001}$ \\
\midrule
1  & 2.706 & 3.841  & 5.024  & 6.635  & 10.828 \\
2  & 4.605 & 5.991  & 7.378  & 9.210  & 13.816 \\
3  & 6.251 & 7.815  & 9.348  & 11.345 & 16.266 \\
4  & 7.779 & 9.488  & 11.143 & 13.277 & 18.467 \\
5  & 9.236 & 11.070 & 12.833 & 15.086 & 20.515 \\
6  & 10.645& 12.592 & 14.449 & 16.812 & 22.458 \\
7  & 12.017& 14.067 & 16.013 & 18.475 & 24.322 \\
8  & 13.362& 15.507 & 17.535 & 20.090 & 26.125 \\
9  & 14.684& 16.919 & 19.023 & 21.666 & 27.877 \\
10 & 15.987& 18.307 & 20.483 & 23.209 & 29.588 \\
\bottomrule
\end{tabular}
\end{center}

\printindex

\end{document}